% Continuous portable body; load g1-appendix.sty in the host.
\GOneBegin
\GOneSection{Production Function Estimation}
\GOneLabel{app:nominal-production}

\GOneSubsection{Introduction}
\GOneLabel{sec:purpose}
We compare the value manufacturers report creating after accounting for the workers and capital they use. The methods support different comparisons: some cover more industry groups, while a comparison shared by all four uses a much smaller set of observations. We therefore examine the estimated input contributions alongside the records on which the methods can be compared.

We take the supplied value-added field as the production measure. Value added represents the value created through production after allowing for purchased inputs; here we retain the reported amount rather than reconstructing it from other reported totals. We call the part left after accounting for labor and capital the \emph{production residual}. A larger residual means more reported value added relative to the input contributions assigned by the estimated equation. Prices, measurement error and temporary production shocks can also enter that remainder, so it does not isolate physical productivity.

We use a modified copy of the supplied annual manufacturing records, called G0.5. Its values have been altered for development, and its company identifiers link records across years through fictional assignments within each province and sector. These assignments allow us to run calculations that require consecutive observations; they do not establish actual company histories. We retain the reported monetary amounts without adjusting them to common prices. Differences in prices and reporting units can therefore affect the results alongside differences in production. We use these records to examine how the calculations behave, rather than to draw empirical conclusions about the productivity of actual manufacturers.

Estimating the input contributions requires more than comparing larger and smaller manufacturers. A company may hire more workers or buy more materials when it expects favorable production conditions that we cannot observe. The observed relation between inputs and production then reflects both the input contribution and the conditions that prompted the company to use them. The four methods we examine address this problem through assumptions about input choices and their timing:
\begin{itemize}
\item Materials purchases can reveal production conditions known to the company, including when it makes no new capital investment. The \textbf{Levinsohn--Petrin (LP)} method uses materials to represent that information when estimating the input contributions.
\item Investment can also convey this information, while differences in the chances of continuing to operate can affect which companies we observe. The \textbf{Olley--Pakes (OP)} method addresses both issues. It requires usable investment observations and information about continuation.
\item Separating labor's contribution from the information conveyed by materials requires additional care about when the company chooses its inputs. The \textbf{Ackerberg--Caves--Frazer (ACF)} method postpones that separation to a later estimation step, using explicit timing assumptions.
\item A system of equations can link current and previous observations to estimate the input contributions together. We use a \textbf{Wooldridge} formulation and specify its version because those equations and their assumptions determine how to read the estimates.
\end{itemize}

Using the observations required by each method, we obtain 37 main sector estimates: LP in 12 sectors, ACF in 12, Wooldridge in 12 and OP in 1. These estimates span 12 manufacturing groups. Requiring inputs suitable for all four methods gives 4 separate common-sample estimates. We compare the methods on shared records and distinguish that comparison from the coverage of their broader samples. LP and Wooldridge place the shared records in closely aligned order, while ACF changes their ordering more. We examine those comparisons together with coefficient uncertainty and sensitivity to estimation choices.

\GOneSubsubsection{Guide to the appendix}
\GOneLabel{sec:roadmap}
\begin{itemize}
\item \textbf{We define what the residual measures.} Section~\GOneRef{sec:meaning} explains how we account for workers and capital, why we use logarithms, and how we express a record's residual relative to its sector average. These steps determine what a high or low reported value means.
\item \textbf{We establish which records supply the calculation.} Section~\GOneRef{sec:data} identifies the source variables, years and missing values. The production observations cover 2011--2015, 2018--2022; the preceding calendar year supplies capital at the start of each year under our year-end-stock assumption. We distinguish the records used in estimation from the broader set for which we can calculate a residual.
\item \textbf{We develop each estimation method.} Section~\GOneRef{sec:methods} explains the common assumptions about input choices. Sections~\GOneRef{sec:lp}, \GOneRef{sec:op}, \GOneRef{sec:acf} and~\GOneRef{sec:wrdg} then present each method's data requirements and equations, followed by the available estimates, uncertainty, distributions and comparisons over time.
\item \textbf{We compare results on the same records.} Section~\GOneRef{sec:coefficients} brings the input coefficients and their uncertainty together. Sections~\GOneRef{sec:distributions} and~\GOneRef{sec:agreement} compare residual distributions and rankings, distinguishing each method's sample from the shared observations used for direct comparison.
\item \textbf{We assess how the estimation choices affect the findings.} Section~\GOneRef{sec:sensitivity} changes the starting values used by the numerical search and the flexibility of the fitted relationships. Section~\GOneRef{sec:synthesis} draws together the findings and their implications for interpreting the residuals.
\end{itemize}

\GOneSubsection{Production residuals}
\GOneLabel{sec:meaning}
\GOneSubsubsection{Input elasticities}
\GOneLabel{sec:elasticities}
We begin by specifying how the measured inputs enter the production equation. Let $i$ identify a company and $t$ its calendar year. In the present data, $i$ refers to the fictional company assignment. Write $Y_{it}$ for reported value added, $L_{it}$ for the number of workers and $K_{it}$ for capital available at the beginning of that year. We assume that these quantities enter the following relation:
\begin{equation}
Y_{it}=A_{it}L_{it}^{\beta_L}K_{it}^{\beta_K}\exp(\varepsilon_{it}),
\GOneLabel{eq:levels}
\end{equation}
The factor $A_{it}$ represents production conditions beyond the two measured inputs. The term $\exp(\varepsilon_{it})$ allows a temporary discrepancy; $\exp$ denotes the exponential function. The parameters $\beta_L$ and $\beta_K$ determine how strongly value added responds to labor and capital. They are \emph{input elasticities}: a small proportional increase in an input changes value added by that proportion multiplied by its elasticity, holding the other terms constant. We estimate these parameters from the records. The multiplicative form in Equation~\GOneEqref{eq:levels} is an assumption we impose, rather than a relationship established by the source files.

Value added per worker provides a simpler comparison, but it leaves differences in capital use mixed with differences in production performance. Equation~\GOneEqref{eq:levels} lets us account for labor and capital separately. Its usefulness depends on how well we estimate their contributions, given that companies can choose inputs in response to production conditions we do not observe.

We work with natural logarithms, written $\ln$, so that multiplying input quantities becomes adding their contributions. This also puts proportional differences on a common scale: the difference between two logged values equals the log of their ratio. Lower-case letters denote the logged quantities:
\begin{equation}
y_{it}=\beta_Ll_{it}+\beta_Kk_{it}+a_{it}+\varepsilon_{it},\qquad
(y,l,k,a)=(\ln Y,\ln L,\ln K,\ln A).
\GOneLabel{eq:logs-first-principles}
\end{equation}
Each term now represents a separate contribution to logged value added. We can take logs only of positive quantities. We therefore count zeros, negative values and unavailable observations separately in the data section and exclude them where the calculation requires a log.

\GOneSubsubsection{Uncentered residual}
\GOneLabel{sec:native}
After estimating the labor and capital coefficients, we multiply each coefficient by its corresponding logged input and subtract both contributions from logged value added. The hat on a coefficient denotes an estimate. We allow the estimates to differ by method $m$ and sector $s$:
\begin{equation}
\widehat r_{it}^{(m)}=y_{it}-\widehat\beta_{L,s}^{(m)}l_{it}-\widehat\beta_{K,s}^{(m)}k_{it}.
\GOneLabel{eq:native-definition}
\end{equation}
This is the \emph{uncentered residual}: we have not yet subtracted a reference average. Under Equation~\GOneEqref{eq:logs-first-principles}, it contains both $a_{it}$ and the temporary discrepancy $\varepsilon_{it}$, as well as any constant that we have not subtracted. Because we use monetary observations at reported prices, the residual can also reflect prices and reporting differences.

A comparison of two records clarifies the interpretation. If they have identical measured labor and capital and use the same fitted coefficient vector, their residual difference equals the difference in logged value added. If their inputs differ, the residual difference subtracts the model-assigned input contributions first. Thus a higher residual means higher reported value added \emph{relative to those contributions}. It does not by itself explain why the difference arose.

The units affect the residual's numerical level. To see why, suppose we multiply value added, labor and capital by constants that change their units: $Y^*=c_YY$, $L^*=c_LL$ and $K^*=c_KK$. Keeping the same estimated coefficients gives
\begin{equation}
\widehat r^*_{it}=\widehat r_{it}+\ln c_Y-\widehat\beta_L\ln c_L-\widehat\beta_K\ln c_K.
\GOneLabel{eq:unit-rescaling}
\end{equation}
A uniform scale change shifts the whole residual series. It need not change the ranking within the same sector, but it makes an unqualified comparison of raw residual levels misleading. Firm-specific price or unit differences are more serious: they need not reduce to a constant shift and can change the ordering itself.

\GOneSubsubsection{Centering}
\GOneLabel{sec:centered}
We use the records in each sector's consecutive-year estimation stage, pooled over all supported years, as the reference. These records require eligible information in the immediately preceding year; the data and method sections give the full selection rules. Let $\mathcal D_s^{(m)}$ contain the company-year observations in this sample for method $m$, and let $N_s^{(m)}$ be their number. Define
\begin{equation}
\overline r_s^{(m)}=\frac{1}{N_s^{(m)}}\sum_{(j,\tau)\in\mathcal D_s^{(m)}}\widehat r_{j\tau}^{(m)},\qquad
\widetilde r_{it}^{(m)}=\widehat r_{it}^{(m)}-\overline r_s^{(m)}.
\GOneLabel{eq:centered-definition}
\end{equation}
The mean centered residual is zero over this pooled reference sample within each sector and method. Centering removes the arbitrary common shift in Equation~\GOneEqref{eq:unit-rescaling} when the scaling constants apply uniformly within the reference group. It does not remove differing prices, product mix, reporting errors or model misspecification. A common zero mean also need not imply equal medians, equal tails or agreement about individual records.

We subtract this same reference mean in every prediction year. We do not recenter each year separately: doing so would remove the annual mean differences we want to describe. Annual means can therefore differ from zero, including in years that contribute to the pooled fit. Their differences describe changing reported outcomes and inputs evaluated under one sector equation, together with changes in the records present. They are neither separate annual coefficient estimates nor estimates of real productivity growth.

\GOneSubsubsection{Relative index}
\GOneLabel{sec:exponentiation}
We can also express the centered log residual as a positive index:
\begin{equation}
I_{it}^{(m)}=\exp\!\left(\widetilde r_{it}^{(m)}\right).
\GOneLabel{eq:relative-index}
\end{equation}
Exponentiation reverses the log transformation. An index above one means that the record's uncentered residual exceeds its sector's reference mean for that method; an index below one means it falls below. Within the same sector and method, the ratio of two indices equals the ratio of their exponentiated uncentered residuals: both use the same reference mean. For example, an index ratio of two means a twofold difference in that model-defined component. It does not imply twice the physical productivity, value added or value added per worker.

An average calculated after taking logs differs from an ordinary average of the index values. Taking the exponential of the average log index gives its geometric mean, which is one in the reference sample because the average centered residual is zero. The arithmetic mean, calculated by adding the index values and dividing by their number, need not be one. Large positive residuals become especially large index values after exponentiation. We consequently report raw logs, centered logs and the positive index separately, with the transformation and denominator stated for each exhibit. These definitions establish what the subsequent distributions measure; we next explain which source observations permit the calculation.

\GOneSubsection{Data and sample construction}
\GOneLabel{sec:data}
We now identify which source observations provide the quantities in Section~\GOneRef{sec:meaning}. G0.5 contains 33,049 records from twelve supplied annual files. We use its modified values and fictional company links as described in the introduction. Each record retains its source year. A shared variable name helps us locate a field across files, but does not establish that the files use the same units or definition.

We use monetary amounts at reported prices, referred to as \emph{nominal} values. Converting them to values at common prices would require suitable price adjustments, or deflators, for value added, capital and materials. The available documentation does not establish a consistent set of those adjustments and source units. We therefore retain each reported scale and treat these fields as provisional measures. Table~\GOneRef{tab:measurement} identifies the quantity and transformation assigned to each variable.
\GOneInput{tables/measurement.tex}
We treat each year's positive reported capital total as its year-end stock. This maintained assumption extends the calculation beyond the years with explicit stock documentation; it does not provide new evidence that the source definitions agree. Beginning capital for production in year $t$ is the reported stock in calendar year $t-1$, linked through the same fictional company assignment. Reported additions separately supply the assumed investment measure for OP. We do not reconstruct investment by subtracting successive stocks, because that difference can also reflect disposals, depreciation, valuation and reporting changes.

The source years are 2010--2015, 2017--2022. Beginning capital permits production observations in 2011--2015, 2018--2022, while the stage requiring eligible production in two consecutive years can use endpoints in 2012--2015, 2019--2022. The unavailable 2016 file supplies neither a stock nor a production observation. In particular, a 2017 production record lacks beginning capital, whereas a 2018 record can use the 2017 stock. A dynamic pair ending in 2018 still lacks an eligible 2017 predecessor. These calendar consequences follow from the timing rule rather than a decision to compress the gap.

We assume a consistent positive unit scale for each variable within a sector over the retained years. This does not correct a reporting-unit error affecting a particular company or year. We pool all supported calendar pairs within each sector with constant production slopes and constant proxy and transition functions, without additional year effects. Pooling therefore draws information from several years by imposing a common relationship across them. Inflation, source changes and changing composition can affect that relationship and the annual residual comparisons.

We estimate each manufacturing group separately, using the two-digit industry codes supplied in the records. This allows the input coefficients to differ across industries. We refer to these groups as sectors; the common digits do not by themselves verify a common classification across editions.

The methods use current and previous observations in different steps. An initial regression can use records with all required positive inputs. A later step relates current production to the same fictional company's preceding year and requires a consecutive pair, which we call a \emph{transition}. We call these the first-stage and dynamic-stage samples, respectively. One company can contribute several transitions; a count of transitions is therefore not a count of distinct companies. The method sections specify the exact calculations, including Wooldridge's system of equations.

We set minimum sample sizes before estimation: 100 observations for the first stage, 50 eligible transitions and 50 distinct fictional companies in the dynamic stage. OP also estimates continuation. For that calculation, a \emph{risk interval} is a pair of calendar years for which we can use the earlier record to determine whether the same fictional assignment appears again. We require 100 usable risk intervals, including 10 disappearances and 10 continuations. These outcomes describe the generated links, not verified company closures. Table~\GOneRef{tab:support} shows how many observations remain under each set of requirements.
\GOneInput{tables/support.tex}
OP has 1 accepted sector estimate and 128 dynamic company-year observations; LP has 12 accepted sector estimates and 5,041 dynamic company-year observations; ACF has 12 accepted sector estimates and 5,041 dynamic company-year observations; Wooldridge has 12 accepted sector estimates and 5,041 dynamic company-year observations. These counts describe each method's own sample. Repeated annual observations from a fictional company remain distinct records. We keep failures of the sample and numerical requirements visible rather than borrowing coefficients from another sector.

Figure~\GOneRef{fig:sample-support} compares these sample sizes. Each method's own sample requires the observations needed by that method. Requiring observations suitable for all four methods imposes additional restrictions; Table~\GOneRef{tab:support} reports that smaller common sample separately.
\begin{figure}[H]
\centering
\caption{Consecutive-year estimation samples}
\GOneLabel{fig:sample-support}
\includegraphics[width=\linewidth]{figures/sample_support.pdf}
\appnotes{Counts pool all supported calendar transitions under each method's own input requirements. Common-input counts appear separately in Table~\GOneRef{tab:support}. Sample thresholds precede numerical estimation; positive support alone does not guarantee a fitted coefficient.}{G0.5 v3 modified manufacturing records with fictional company links; authors' calculations.}
\end{figure}
A method's own sample uses its required inputs; the common sample requires inputs suitable for all four methods before applying the same calendar rule. We report accepted and rejected estimation cells under both definitions. Shared input eligibility alone does not guarantee that every estimator passes its numerical and selection requirements.

\GOneSubsubsection{Input availability and response states}
\GOneLabel{sec:deep-inputs}
We first count how many records provide each required variable in every supplied year. Table~\GOneRef{tab:deep-input-availability} reports these counts separately, so that missing capital does not obscure the availability of value added or workers. Estimation requires the necessary variables together in the same record, and the dynamic stage uses consecutive-year information, requiring an eligible record in the preceding year as well. The separate counts therefore describe the available information before we form the estimation sample.
\GOneInput{tables/deep_input_availability.tex}
We extend beginning capital to every available year with a positive stock in its immediately preceding calendar year. This uses the stated year-end assumption even where documentary support is narrower. The missing 2016 file leaves no 2017 beginning stock. A 2018 observation can supply current inputs using the 2017 stock, but a dynamic pair ending in 2018 also needs an eligible 2017 production observation. The source gap therefore affects prediction and dynamic estimation differently. We report those sample differences explicitly.

To explain why a value cannot enter the calculation, we distinguish a missing response from a reported zero or negative amount. Table~\GOneRef{tab:deep-input-states} counts these response states for the history used in the model. We also report failures of the maintained assumption-domain and timing rules. The stock and investment meanings are maintained assumptions for every supplied year; we retain the documentary evidence without treating it as verification of those assumptions.
\GOneInput{tables/deep_input_states.tex}
We cannot take the natural logarithm of zero. Adding a small constant would change the measured proxy and the records used to estimate its relationship with output, so we keep the specified positive-value rule. Missing beginning capital has a different meaning: the prior stock may be unavailable even if the record reports current production. Neither exclusion establishes that an actual company used no capital.

We next examine how much the positive amounts vary before requiring complete inputs on the same record. Table~\GOneRef{tab:deep-input-levels} reports the median and selected percentiles in the supplied units. The median divides the ordered observations in half; P25 and P75 locate the lower and upper quarters, while P1 and P99 describe values near the ends. Monetary amounts and worker counts retain different units even when we later take their logarithms.
\GOneInput{tables/deep_input_levels.tex}
The upper percentiles are well above the median for several inputs, showing that a typical positive observation does not describe the largest reported amounts. This variation can affect the estimated relationship between inputs and output as well as the residual distribution. A large amount alone does not establish a reporting error. We retain eligible observations, including the large values, and examine the resulting residual tails below.

\GOneSubsubsection{Sample restrictions}
\GOneLabel{sec:deep-funnel}
We now require the inputs together and then apply the calendar condition. We first retain records with all required current inputs; the dynamic sample also requires an eligible predecessor. Table~\GOneRef{tab:deep-funnel} follows these restrictions in order for the materials-based methods. At each step, we count exclusions among the records retained by the preceding step. We can therefore distinguish a missing or unusable beginning stock from the later loss of a record that lacks an eligible predecessor.
\GOneInput{tables/deep_funnel.tex}
The table separates losses due to missing beginning capital from losses due to a missing eligible predecessor. We apply the minimum estimation requirements after pooling the valid pairs within each sector. A broader set can receive residuals because calculating a residual only requires current inputs and an available coefficient vector. Records outside the dynamic stage can also contribute to an earlier regression in methods with a separate first stage. We therefore retain each record's role rather than treating every prediction as an estimation observation.
Figure~\GOneRef{fig:deep-input-core} compares the number of positive values for each input with the full annual record count. We draw a separate bar for each variable because an observation can supply one input and lack another. Estimation still requires the inputs together and the specified calendar history.
\begin{figure}[H]
\centering
\caption{Availability of outcome, labor and capital}
\GOneLabel{fig:deep-input-core}
\includegraphics[width=\linewidth]{figures/deep_input_availability_core.pdf}
\appnotes{Bars count positive reported or constructed values separately for each variable and source year; the background shows the full annual sample. Positive values still require the applicable calendar timing before estimation; stock and investment meanings follow the stated assumptions. The legend distinguishes the missing annual file from reported zero counts. We retain zero counts. These separate-variable counts do not give the number of records used to estimate coefficients.}{G0.5 v3 modified manufacturing records with fictional company links; authors' calculations.}
\end{figure}
Positive amounts can still fail to provide all required inputs for the same record. Beginning capital also needs the immediately preceding stock, and dynamic estimation needs an eligible preceding production observation. The annual bars show the input availability before imposing those joint requirements.
Figure~\GOneRef{fig:deep-input-proxy} compares the number of positive values for each input with the full annual record count. We draw a separate bar for each variable because an observation can supply one input and lack another. Estimation still requires the inputs together and the specified calendar history.
\begin{figure}[H]
\centering
\caption{Availability of proxy inputs}
\GOneLabel{fig:deep-input-proxy}
\includegraphics[width=\linewidth]{figures/deep_input_availability_proxy.pdf}
\appnotes{Bars count positive reported or constructed values separately for each variable and source year; the background shows the full annual sample. Positive values still require the applicable calendar timing before estimation; stock and investment meanings follow the stated assumptions. The legend distinguishes the missing annual file from reported zero counts. We retain zero counts. These separate-variable counts do not give the number of records used to estimate coefficients.}{G0.5 v3 modified manufacturing records with fictional company links; authors' calculations.}
\end{figure}
Positive amounts can still fail to provide all required inputs for the same record. Beginning capital also needs the immediately preceding stock, and dynamic estimation needs an eligible preceding production observation. The annual bars show the input availability before imposing those joint requirements.


\GOneSubsection{Proxy-variable estimators}
\GOneLabel{sec:methods}
We want to compare reported value added after accounting for how much labor and capital each company uses. Equation~\GOneEqref{eq:native-definition} gives this calculation once we know the coefficients that measure the two inputs' contributions. The difficulty is estimating those coefficients. A company that expects favorable production conditions may hire more workers or purchase more materials. A regression can then attribute some of the favorable conditions' contribution to the inputs, because both contributed to the higher output. Econometricians call this the simultaneity problem: input choices and the unexplained part of production depend on the same information.

We separate the unexplained part of logged value added into conditions the company knows when choosing its inputs, $\omega_{it}$, and a remaining discrepancy that it does not know at that point, $\varepsilon_{it}$. With $y$, $l$ and $k$ denoting the logarithms of reported value added, workers and beginning-of-year capital for company $i$ in year $t$, the resulting equation is
\begin{equation}
y_{it}=\beta_Ll_{it}+\beta_Kk_{it}+\omega_{it}+\varepsilon_{it}.
\GOneLabel{eq:production}
\end{equation}
The coefficients $\beta_L$ and $\beta_K$ measure the contributions assigned to labor and capital. Ordinary least squares chooses these coefficients to minimize squared prediction errors. Interpreting them as the inputs' separate economic contributions requires, among other conditions, that the inputs be unrelated to the unexplained term. The company's response to $\omega$ is a reason that condition can fail.

We therefore examine methods that use an observed company decision to learn about production conditions we cannot observe directly. Such a decision serves as a \emph{proxy}: materials purchases or investment may reveal whether the company expects favorable conditions. Levinsohn--Petrin (LP) uses materials; Olley--Pakes (OP) uses investment and also addresses which companies remain in the sample. Ackerberg--Caves--Frazer (ACF) revises how we separate labor's contribution from the information in materials purchases. Wooldridge combines equations based on current and previous-year information. Each approach requires assumptions about company decisions and how production conditions change over time. We first explain those common assumptions, then give the equations and results for each method.

\GOneSubsubsection{Materials demand and input timing}
\GOneLabel{sec:proxy-inversion}
Suppose two companies have the same capital and face the same other determinants of materials demand. If the company with more favorable production conditions always buys more materials, we can work backward from its purchases to the conditions that motivated them. Formally, write materials demand as $M_{it}=\mathcal M_t(K_{it},\omega_{it})$ and require it to increase strictly with $\omega$ at a given $K$. This requirement is called strict monotonicity. Working backward through the demand rule, or \emph{inverting} it, gives the following expression in logarithms:
\begin{equation}
\omega_{it}=h_t(k_{it},m_{it}).
\GOneLabel{eq:proxy-inversion}
\end{equation}
The function $h_t$ associates each combination of capital and materials with one value of $\omega$. This one-to-one relationship requires more than positive purchases. Positive values allow us to calculate logarithms; they do not show that a larger materials purchase always reflects more favorable production conditions. Differences in input prices, customer demand, financing constraints or other unobserved shocks may also affect purchases. We use the available nominal materials measure to implement the calculation, while retaining these limits on its economic interpretation. Within each sector, we impose the same control function over all pooled years, replacing the general time-varying function by one common polynomial. We likewise use one transition relationship across the supported calendar pairs and add no year effects. Common price movements and changes in the composition of the samples can therefore enter the estimated relation.

We use materials to account for $\omega$ through Equation~\GOneEqref{eq:proxy-inversion}; this role gives $h_t$ its name, the \emph{control function}. Equation~\GOneEqref{eq:production} therefore has coefficients for labor and capital, while materials enter through that function. Estimating a separate materials contribution would require a different production equation and a fresh assessment of the outcome and input timing. In particular, a gross-output equation would answer a different accounting question from the value-added equation used here.

We also use the previous year's production conditions to predict the current year's conditions. Companies may retain advantages or disadvantages from one year to the next. We represent that predictable part by $g(\omega_{i,t-1})$ and call the remaining new component $\xi_{it}$ the \emph{productivity innovation}:
\begin{equation}
\omega_{it}=g(\omega_{i,t-1})+\xi_{it},\qquad
E[\xi_{it}\mid\mathcal I_{i,t-1},k_{it}]=0.
\GOneLabel{eq:markov-transition}
\end{equation}
The expectation restriction says that the new component has mean zero after we use the relevant past information, $\mathcal I_{i,t-1}$, and beginning-of-year capital, $k_{it}$. Capital already in place at the start of the year cannot be changed in response to an unexpected event later that year. Under this timing assumption, we call it \emph{predetermined} relative to the current innovation. Current employment and materials purchases may respond to the event. Previous-year employment and beginning-of-year capital can therefore help estimate the input coefficients without using current choices that respond to the same surprise. Whether companies actually make their decisions in this order is an economic assumption. The available columns, particularly when linked through fictional company IDs, do not establish it.

\GOneSubsubsection{Polynomial approximation}
\GOneLabel{sec:common-basis}
The functions relating materials to production conditions, and past conditions to current conditions, are unknown. We approximate them with polynomial regressions, which let a fitted relationship bend and let its slope depend on more than one variable. For capital $k$ and a logged proxy $p$, we include their levels, squares, cubes and all products whose total degree is at most three. We collect these terms in the vector
\begin{equation}
B_3(k,p)=(k,p,k^2,kp,p^2,k^3,k^2p,kp^2,p^3)' .
\GOneLabel{eq:polynomial-basis}
\end{equation}
We estimate an intercept separately. The ACF first-stage regression adds labor and the corresponding powers and products of labor, capital and materials. Our primary estimates use the cubic polynomial above. To examine how the approximation affects the estimates, we also fit a quadratic polynomial, which drops terms of degree three. LP, OP and ACF keep a separate cubic relationship between current and previous-year productivity in both specifications. For Wooldridge, we change the current and previous-year proxy polynomials together, as explained in Section~\GOneRef{sec:wrdg}.

Before estimating each sector and method, we subtract each logged variable's median in that method's eligible source sample. We subtract the same median from its previous-year values. This makes the numbers used in the polynomial easier to handle numerically. Because we include every polynomial term up to the chosen degree, subtracting a constant changes how we write the fitted function without changing the set of functions the polynomial can represent. A numerical search can nevertheless respond to the changed scale and location of its inputs, which is why we examine convergence and sensitivity to starting values.

After estimation, we calculate residuals using the estimated labor and capital coefficients and the original log inputs in Equation~\GOneEqref{eq:native-definition}. Subtracting input medians before estimation therefore serves a different purpose from subtracting sector residual means afterward. The latter calculation, in Equation~\GOneEqref{eq:centered-definition}, gives each sector a common residual reference level for comparison.

\GOneSubsubsection{Coefficient uncertainty}
\GOneLabel{sec:method-uncertainty}
An estimated coefficient depends on the observations used to calculate it. We examine this source of uncertainty by drawing repeated samples of fictional companies, with replacement, and re-estimating the model in each sample. We resample each company's complete history together so that its consecutive-year records remain linked, and we repeat every estimation stage. This procedure is the company-history bootstrap.

We report the standard deviation of the resulting coefficient estimates as the bootstrap \emph{standard error}. A larger value means that the estimate varies more across the resampled histories. We also report a percentile interval covering the middle 95\% of those estimates, which shows where most of the resampled coefficients lie without imposing a symmetric interval around the original estimate. These intervals are exploratory and conditional on the fictional links, measurement choices and estimator. They do not measure uncertainty about whether a fictional match represents a real company, nor do they incorporate unresolved prices or source definitions. Numerical convergence in every resample also does not establish that a nonlinear search has found a unique solution.

We begin with Levinsohn--Petrin because its two stages show how materials purchases and previous-year information play different roles. We then examine what investment and sample-continuation information would add under Olley--Pakes, before turning to Ackerberg--Caves--Frazer and Wooldridge. For each method, we specify the calculation we actually use; the method's name alone does not determine its polynomial, input timing or estimation criterion.

\GOneSubsection{Levinsohn--Petrin}
\GOneLabel{sec:lp}
\GOneSubsubsection{Materials proxy}
\GOneLabel{sec:lp-economics}
Levinsohn and Petrin use a company's purchases of intermediate inputs, such as materials, to account for production conditions that the company observes and the researcher does not \citep{lp}. A company may purchase materials every year even if it invests in new capital only occasionally. Materials can therefore supply information in years with no positive investment. The separate availability of materials and investment is therefore important when comparing the methods' samples. Materials purchases can reveal unobserved conditions because companies respond to those conditions; using purchases this way still requires the demand assumptions described above.

\GOneSubsubsection{Variables and timing}
\GOneLabel{sec:lp-data}
We use the logarithms of positive reported nominal value added for $y$, positive worker counts for $l$, and positive reported raw/supporting materials for $m$. For $k$, we use the logarithm of the capital stock reported for the preceding December, under the maintained year-end-stock assumption. Production in year $t$ uses the positive reported stock in calendar year $t-1$. The first-stage regression pools all observations with eligible current inputs. The second-stage calculation requires an eligible predecessor linked by the same fictional company ID, and pools every supported calendar pair within the sector.

The two stages consequently use different sample sizes, as Table~\GOneRef{tab:support} shows. An eligible production record can inform the first-stage regression even if it has no eligible predecessor or continuation. The second stage also needs the previous-year capital and materials values; usable current values cannot fill those gaps. We check each stage's minimum sample requirement before estimation. When resampling company histories to calculate uncertainty, we reapply these same sample rules.

\GOneDetail{Source variables}
\GOneLabel{sec:lp-source-variables}
We check the fields in Table~\GOneRef{tab:deep-lp-mapping} against the requirements for LP estimation. The table links each variable in the equations to its source field and states the condition a record must meet before we take its logarithm.
\GOneInput{tables/deep_lp_mapping.tex}
The outcome measures reported value added at reported prices. We measure labor by the number of workers, without information about their hours or skills. For capital, we assume that the preceding year's reported stock describes the assets available at the start of the production year. This assumption extends the calculation to the available years without establishing that their source definitions are equivalent. A stock also does not directly measure how intensively a company uses its assets. These choices affect both the estimated input contributions and the remainder we compare.

\GOneDetail{Year coverage}
\GOneLabel{sec:lp-year-coverage}
For production in year $t$, we obtain beginning capital from the stock reported in calendar year $t-1$, treating that amount as its year-end stock. Materials or investment comes from year $t$ itself. The dynamic stage also needs the preceding production observation; its beginning capital consequently draws on year $t-2$. We preserve every calendar gap instead of treating the next observed year as the next calendar year. Table~\GOneRef{tab:deep-lp-input-states} reports each input's availability before we require them together.
\GOneInput{tables/deep_lp_input_states.tex}
The positive observations for one input need not be the same observations as those for another. We retain a record only when it supplies all the quantities that its estimation step requires. A missing response does not show zero production, and we cannot take the log of a recorded zero. Beginning capital also requires the immediately preceding stock. The year-end and investment assumptions admit positive reported amounts without changing the documentary flags; those flags continue to show where the source evidence is narrower than the assumption.

\GOneDetail{Transformations and sample restrictions}
\GOneLabel{sec:lp-transformations}
We check positivity before taking logs. We use the available production years and sectors, match each record to the immediately preceding calendar year, and apply the minimum sample sizes to the pooled sector sample. For the polynomial calculations, we then subtract the relevant sample median from each logged variable; this changes its numerical reference point while keeping the eligible records fixed. We estimate one coefficient vector and control relationship per sector over the pooled years, without adding year effects. Table~\GOneRef{tab:deep-lp-sample} traces the sample by sector and year.
\GOneInput{tables/deep_lp_sample.tex}
The estimation steps use these records in different ways. An earlier observation can supply the preceding-year inputs for a later calendar pair. The estimators with a separate first stage can also use eligible records outside the accepted pairs in that first-stage calculation. Once a sector has estimated coefficients, we calculate a residual for each record with the required current inputs, including one without an eligible predecessor. Calculating that residual does not add a new pair or re-estimate the coefficients. An observation outside the dynamic-stage sample can therefore contribute to the first stage, receive a residual, or do both.


\GOneSubsubsection{Estimation procedure}
\GOneLabel{sec:lp-stages}
Substituting Equation~\GOneEqref{eq:proxy-inversion} into the production equation gives
\begin{equation}
y_{it}=\beta_Ll_{it}+\phi(k_{it},m_{it})+\varepsilon_{it},\qquad
\phi(k,m)=\beta_Kk+h(k,m).
\GOneLabel{eq:lp-first-population}
\end{equation}
The function $\phi$ contains both capital's contribution, $\beta_Kk$, and the unobserved conditions represented by $h$. The equation alone cannot separate them: changing the capital coefficient and offsetting that change inside $h$ would give the same predicted output. We therefore estimate capital's contribution in a second stage, using changes over time. In the first stage, the original LP procedure estimates labor's coefficient separately from $\phi$. For that separation to have an economic interpretation, labor must vary in ways that the capital and materials function cannot already explain. Section~\GOneRef{sec:acf} examines why choosing labor and materials using the same information can undermine that requirement.

We first regress logged value added on logged workers, an intercept and the capital--materials polynomial $B_3(k,m)$, using records with positive values for all required inputs and the outcome. Let $\widehat y^{\,1}$ denote this regression's fitted value. We retain its labor coefficient, $\widehat\beta_L$, and subtract the fitted labor contribution from predicted output. We then try a capital coefficient, $b_K$, and subtract the capital contribution it would imply:
\begin{equation}
\widehat\phi_{it}=\widehat y^{\,1}_{it}-\widehat\beta_Ll_{it},\qquad
\widehat\omega_{it}(b_K)=\widehat\phi_{it}-b_Kk_{it}.
\GOneLabel{eq:lp-phi}
\end{equation}
Each choice of $b_K$ gives a different provisional estimate of $\omega$. We calculate that estimate for both years of every eligible company pair one calendar year apart. We then regress its current value on an intercept and the first three powers of its previous-year value. With $\mathcal D$ denoting the sample of these consecutive-year pairs, the calculation is
\begin{equation}
\begin{aligned}
H_{it}(b_K)&=(1,w_{it},w_{it}^2,w_{it}^3),\\
w_{it}&=\widehat\phi_{i,t-1}-b_Kk_{i,t-1},\\
\widehat\gamma(b_K)&=\arg\min_\gamma\sum_{(i,t)\in\mathcal D}
 [\widehat\omega_{it}(b_K)-H_{it}(b_K)\gamma]^2.
\end{aligned}
\GOneLabel{eq:lp-transition}
\end{equation}
The fitted value from this regression estimates the predictable component in Equation~\GOneEqref{eq:markov-transition}: how much of current productivity we could predict from the previous year. For every trial capital coefficient, we refit this relationship, obtain $\widehat\gamma(b_K)$ and calculate the resulting error in predicted value added. We select the capital coefficient that gives the smallest sum of squared errors:
\begin{equation}
\widehat\beta_K=\arg\min_{b_K}\sum_{(i,t)\in\mathcal D}
 [y_{it}-\widehat\beta_Ll_{it}-b_Kk_{it}-H_{it}(b_K)\widehat\gamma(b_K)]^2.
\GOneLabel{eq:lp-objective}
\end{equation}
This calculation is called \emph{profiled nonlinear least squares}. We search over the capital coefficient while solving for the productivity-transition coefficients at each trial value; those transition coefficients are thus ``profiled out'' of the outer search. Equations~\GOneEqref{eq:lp-phi}--\GOneEqref{eq:lp-objective} define the LP implementation reported here, drawing on the implementation family described by \citet{prodest}. Other implementations use different estimation criteria. Our reported calculation uses this squared-error criterion and does not include a correction for which companies remain in the sample.

\GOneSubsubsection{Worked aggregate calculation}
\GOneLabel{sec:lp-worked}
We illustrate the calculation using mean logged inputs within one sector. Multiplying each mean log input by its fitted coefficient and subtracting both contributions from mean logged value added gives the mean individual residual exactly, because the same slopes apply to every record in that sector. Table~\GOneRef{tab:deep-lp-worked} follows Equations~\GOneEqref{eq:native-definition}--\GOneEqref{eq:relative-index} through centering and exponentiation, using full precision before rounding.
\GOneInput{tables/deep_lp_worked.tex}
In 2011, mean logged value added is 13.7219. The fitted labor and capital contributions are 2.0838 and 1.1658, leaving 10.4723. Subtracting the pooled reference mean of 11.1596 gives a centered mean of -0.6873. Its exponential is 0.5029, while the arithmetic mean of individual index values is 2.1457. The latter gives greater weight to large individual residuals after exponentiation. Both annual examples use the same pooled equation; neither estimates a separate annual technology.


\GOneSubsubsection{Coefficient estimates}
\GOneLabel{sec:lp-coefficients}
We record the primary outcome for every sector and both input-sample definitions in Table~\GOneRef{tab:deep-lp-attempts}. This keeps the absence of a fit separate from a coefficient estimated near zero.
\GOneInput{tables/deep_lp_attempts.tex}
Only estimated cells supply coefficients and derived residuals. A count requirement and a numerical failure have different meanings: the former prevents the declared calculation from being attempted, while the latter prevents an attempted calculation from supplying an accepted estimate.
We report the estimated coefficients beginning with Table~\GOneRef{tab:deep-lp-regression}. The own and common samples have separate columns and fits. To describe uncertainty, we resample whole fictional company histories and repeat every estimation step. Standard errors measure the spread of these repeated estimates; percentile intervals give their admitted lower and upper quantiles.
\GOneInput{tables/deep_lp_regression.tex}
\GOneInput{tables/deep_lp_regression-2.tex}
\GOneInput{tables/deep_lp_regression-3.tex}
\GOneInput{tables/deep_lp_regression-common.tex}
The uncertainty in the coefficient sum depends on how labor and capital estimates move together across resamples. Adding their separate standard errors would ignore that covariance. Fitted year coverage is pooled, and a separate first stage can use more records than the dynamic sample. These counts describe distinct roles in the same estimation.

Table~\GOneRef{tab:deep-lp-numerics} reports the rank and optimization diagnostics for the successful primary fits. We compare numerical rank with the dimension of the relevant matrix because returning a coefficient is not enough to show that the stated calculation had the required independent information.
\GOneInput{tables/deep_lp_numerics.tex}
The diagnostics describe the numerical problem at the reported solution. They do not establish valid economic timing or a unique global optimum; the alternative-start comparison examines a separate aspect of that uncertainty.

Figure~\GOneRef{fig:deep-lp-coefficients} displays the same own-sample coefficients and full admitted uncertainty intervals. Separate input panels preserve the scale of the labor and capital estimates.
\begin{figure}[H]
\centering
\caption{LP: coefficient estimates}
\GOneLabel{fig:deep-lp-coefficients}
\includegraphics[width=\linewidth]{figures/deep_lp_coefficients.pdf}
\appnotes{We pool supported years within each sector under constant coefficients. Marks show own-sample primary estimates and intervals admitted under the bootstrap minimum-success rule. Intervals retain their full endpoints and condition on the fictional histories and maintained measurements.}{G0.5 v3 modified manufacturing records with fictional company links; authors' calculations.}
\end{figure}
Across the estimated sectors, labor ranges from 0.1842 to 0.6574; capital ranges from -0.0209 to 0.0867. We retain the signs and magnitudes returned by the primary specification. Their interpretation depends on beginning-capital timing, reported monetary scales, the proxy assumptions and the estimation criterion. The intervals describe additional variation when the same calculation is repeated on resampled histories. The capital slope is negative in sectors 14 and 22. At given reported value added and workers, higher beginning capital therefore raises the calculated residual in those equations. This sign is a feature of the nominal fitted relationship; it does not establish that additional productive capital reduces actual output. It changes how capital differences enter the residual ranking and must be considered when interpreting that ordering.
\GOneSubsubsection{Bootstrap coefficient distributions}
\GOneLabel{sec:lp-bootstrap}
Table~\GOneRef{tab:deep-lp-bootstrap} reports the spread of the repeated coefficient estimates. The standard error is their sample standard deviation, while the median and percentiles reveal features that a single standard error can conceal. The minimum and maximum make finite extremes visible.
\GOneInput{tables/deep_lp_bootstrap.tex}
The largest labor-coefficient standard deviation occurs in sector 12, at 0.0743. Its bootstrap median is 0.1726, with P5 and P95 at 0.0574 and 0.3118; the full range is 0.0180 to 0.3624. Comparing the middle of this distribution with its extremes helps distinguish broad dispersion from a small number of distant estimates. We retain every finite resampled estimate and keep the original-sample coefficient as the primary result. A limited number of resamples determines the extreme percentiles, so their locations should not be interpreted as precise tail probabilities. The exercise also leaves uncertainty about actual identities, prices and economic restrictions outside the resampling calculation.
\GOneSubsubsection{Residual distributions}
\GOneLabel{sec:lp-distribution}
We subtract the fitted labor and capital contributions from logged value added, then subtract each sector's pooled dynamic residual mean. Figure~\GOneRef{fig:deep-lp-histogram} compares the uncentered and centered distributions on exactly the same fitted records.
\begin{figure}[H]
\centering
\caption{LP: pooled residual distributions}
\GOneLabel{fig:deep-lp-histogram}
\includegraphics[width=\linewidth]{figures/deep_lp_histogram.pdf}
\appnotes{Both panels contain the same own-sample dynamic observations pooled across all fitted years. Uncentered residuals use original logs; centering subtracts one mean per sector-method-sample. Bins span 0.5 log points and retain all finite tails. Labels show the modal and first and last occupied bins with source-derived counts and shares.}{G0.5 v3 modified manufacturing records with fictional company links; authors' calculations.}
\end{figure}
The centered distribution has a standard deviation of 1.467 log points. Centering removes differences in sector means while preserving all within-sector gaps and rankings. The pooled shape can change because it previously combined those mean differences with variation within sectors. We retain all finite observations, including the tails.
\GOneSubsubsection{Prediction samples and tails}
\GOneLabel{sec:lp-prediction}
The histogram above describes the dynamic estimation sample. We can calculate residuals more broadly whenever a record supplies the required current inputs in a sector with estimated slopes. Table~\GOneRef{tab:deep-lp-scopes} distinguishes these roles for every source year and for the pooled records.
\GOneInput{tables/deep_lp_scopes.tex}
There are 10,315 predicted company-year observations: 5,041 belong to the pooled dynamic fit and 5,274 lie outside it. Only the dynamic-fit records determine the reference mean. The annual fit means and the broader prediction means can therefore differ from zero, even though the pooled fit mean is zero by construction. A company appearing repeatedly contributes several observations to these counts.

Table~\GOneRef{tab:deep-lp-quantiles} holds the fitted records fixed and compares the uncentered logs, centered logs and positive index. We include the minimum and maximum so that the transformed tails remain visible.
\GOneInput{tables/deep_lp_quantiles.tex}
Centering preserves every within-sector gap. Exponentiation preserves within-sector ordering but enlarges distances at high values, as Equation~\GOneEqref{eq:relative-index} shows. A large index maximum therefore needs to be read through this transformation, rather than treated directly as physical efficiency.

Table~\GOneRef{tab:deep-lp-tails} identifies distant observations using two distinct rules. The interquartile range (IQR) is the distance from P25 to P75. Fences placed 1.5 IQR beyond those quartiles compare a value with the spread of its own distribution. Fixed absolute-log and index cutoffs instead apply the same numerical threshold in every sector. A record can exceed more than one cutoff.
\GOneInput{tables/deep_lp_tails.tex}
In the pooled dynamic sample, 743 records have an absolute centered log residual above two. These flags describe the residuals produced by the fitted equation. Prices, reported amounts, the estimated input contributions and production conditions can all affect an extreme residual, so the flag alone identifies neither an error nor its cause. We retain every flagged observation.

\GOneSubsubsection{Annual distributions}
\GOneLabel{sec:lp-annual}
Table~\GOneRef{tab:deep-lp-annual} compares all eligible predictions by sector and year. We keep each sector's fitted slopes and reference mean constant over time. This preserves annual differences instead of subtracting each year's average separately.
\GOneInput{tables/deep_lp_annual.tex}
Changes in annual means and quartiles combine differences in reported outcomes and inputs with differences in sample composition. A year with more high-residual observations can have a higher mean without every continuing assignment improving. We examine this distinction next by matching observations within adjacent years.
Figure~\GOneRef{fig:deep-lp-year-sector} and its continuation panels plot the annual median and quartiles for each sector. These intervals measure dispersion across observations, not uncertainty in an estimated coefficient.
\begin{figure}[H]
\centering
\caption{LP annual residual quartiles}
\GOneLabel{fig:deep-lp-year-sector}
\includegraphics[width=\linewidth]{figures/deep_lp_year_sector.pdf}
\appnotes{Each sector panel includes the available annual predictions, using one pooled coefficient vector and reference mean. Dots mark medians and intervals run from P25 to P75; all available values are numerically labelled. An empty group is not a zero. The legend identifies any missing-file region. Continuation panels use the same construction.}{G0.5 v3 modified manufacturing records with fictional company links; authors' calculations.}
\end{figure}
\begin{figure}[H]
\centering
\caption{LP annual residual quartiles (continued)}
\GOneLabel{fig:deep-lp-year-sector-p2}
\includegraphics[width=\linewidth]{figures/deep_lp_year_sector_p2.pdf}
\appnotes{Each sector panel includes the available annual predictions, using one pooled coefficient vector and reference mean. Dots mark medians and intervals run from P25 to P75; all available values are numerically labelled. An empty group is not a zero. The legend identifies any missing-file region. Continuation panels use the same construction.}{G0.5 v3 modified manufacturing records with fictional company links; authors' calculations.}
\end{figure}
\begin{figure}[H]
\centering
\caption{LP annual residual quartiles (continued)}
\GOneLabel{fig:deep-lp-year-sector-p3}
\includegraphics[width=\linewidth]{figures/deep_lp_year_sector_p3.pdf}
\appnotes{Each sector panel includes the available annual predictions, using one pooled coefficient vector and reference mean. Dots mark medians and intervals run from P25 to P75; all available values are numerically labelled. An empty group is not a zero. The legend identifies any missing-file region. Continuation panels use the same construction.}{G0.5 v3 modified manufacturing records with fictional company links; authors' calculations.}
\end{figure}

The sector panels separate variation within each industry from changes in which industries dominate a pooled annual distribution. They describe the same fixed sector equations over time.
\GOneSubsubsection{Adjacent-year changes}
\GOneLabel{sec:lp-paired}
We hold membership fixed within each comparison by requiring a prediction for the same fictional assignment and sector in adjacent calendar years. Table~\GOneRef{tab:deep-lp-pair-coverage} shows how this restriction narrows the endpoint samples. We retain intervals with no matched predictions so that missing capital history remains visible.
\GOneInput{tables/deep_lp_pair_coverage.tex}
The pooled comparison contains 5,041 record-pair intervals. A company can enter several successive intervals, so this is not a count of independent companies. Restricting to matched predictions removes changes in membership within each interval, although membership can still differ between intervals.

For each pair, we subtract the earlier centered residual from the later one. Let $\Delta x_{it}=x_{it}-x_{i,t-1}$ for any logged variable $x$. Holding the sector slopes fixed gives
\begin{equation}
\Delta\widetilde r_{it}=\Delta y_{it}-\widehat\beta_{L,s}\Delta l_{it}-\widehat\beta_{K,s}\Delta k_{it}.
\GOneLabel{eq:lp-paired-change}
\end{equation}
The pooled reference mean cancels because we subtract the same constant in both years. Table~\GOneRef{tab:deep-lp-paired} reports the resulting changes for each interval and sector, then pools the record-pair intervals under the same rule.
\GOneInput{tables/deep_lp_paired.tex}
The pooled mean log change is 0.145. The correlation of earlier and later ranks is 0.793, measuring persistence in ordering. The mean individual index ratio is 1.946; large ratios can strongly influence this arithmetic mean, which need not equal the exponential of the mean log change. These comparisons use fictional links and reported-price equations; they do not establish real-company productivity changes.
Figure~\GOneRef{fig:deep-lp-paired-change} and its continuation panels plot these adjacent-year changes under the same matched-prediction rule.
\begin{figure}[H]
\centering
\caption{LP adjacent-year residual changes}
\GOneLabel{fig:deep-lp-paired-change}
\includegraphics[width=\linewidth]{figures/deep_lp_paired_change.pdf}
\appnotes{Each comparison uses same-sector fictional assignments with eligible predictions in consecutive calendar years. Sector coefficients and pooled reference means remain fixed. Histograms retain all finite changes. Numeric labels identify the modal and first and last occupied bins. Pooled panels combine record-pair intervals, allowing repeated company contributions; interval panels retain their own denominators.}{G0.5 v3 modified manufacturing records with fictional company links; authors' calculations.}
\end{figure}
\begin{figure}[H]
\centering
\caption{LP adjacent-year residual changes (continued)}
\GOneLabel{fig:deep-lp-paired-change-p2}
\includegraphics[width=\linewidth]{figures/deep_lp_paired_change_p2.pdf}
\appnotes{Each comparison uses same-sector fictional assignments with eligible predictions in consecutive calendar years. Sector coefficients and pooled reference means remain fixed. Histograms retain all finite changes. Numeric labels identify the modal and first and last occupied bins. Pooled panels combine record-pair intervals, allowing repeated company contributions; interval panels retain their own denominators.}{G0.5 v3 modified manufacturing records with fictional company links; authors' calculations.}
\end{figure}
\begin{figure}[H]
\centering
\caption{LP adjacent-year residual changes (continued)}
\GOneLabel{fig:deep-lp-paired-change-p3}
\includegraphics[width=\linewidth]{figures/deep_lp_paired_change_p3.pdf}
\appnotes{Each comparison uses same-sector fictional assignments with eligible predictions in consecutive calendar years. Sector coefficients and pooled reference means remain fixed. Histograms retain all finite changes. Numeric labels identify the modal and first and last occupied bins. Pooled panels combine record-pair intervals, allowing repeated company contributions; interval panels retain their own denominators.}{G0.5 v3 modified manufacturing records with fictional company links; authors' calculations.}
\end{figure}
\begin{figure}[H]
\centering
\caption{LP adjacent-year residual changes (continued)}
\GOneLabel{fig:deep-lp-paired-change-p4}
\includegraphics[width=\linewidth]{figures/deep_lp_paired_change_p4.pdf}
\appnotes{Each comparison uses same-sector fictional assignments with eligible predictions in consecutive calendar years. Sector coefficients and pooled reference means remain fixed. Histograms retain all finite changes. Numeric labels identify the modal and first and last occupied bins. Pooled panels combine record-pair intervals, allowing repeated company contributions; interval panels retain their own denominators.}{G0.5 v3 modified manufacturing records with fictional company links; authors' calculations.}
\end{figure}

A pooled average can conceal increases and decreases across matched assignments. The interval panels show how their distributions vary over time under the fixed sector equations.
\GOneSubsubsection{Specification sensitivity}
\GOneLabel{sec:lp-sensitivity}
Figure~\GOneRef{fig:deep-lp-sensitivity} and its continuation panels compare the prespecified starting values and proxy-polynomial alternatives. Changing the start leaves the model fixed; changing the polynomial alters the fitted control relationship. A numerical stopping rule need not identify the same solution from every starting value.
\begin{figure}[H]
\centering
\caption{LP: estimation sensitivity}
\GOneLabel{fig:deep-lp-sensitivity}
\includegraphics[width=\linewidth]{figures/deep_lp_sensitivity.pdf}
\appnotes{The primary cubic specification, two alternative starts and quadratic proxy use the same own-sample eligibility and numerical translations. Objective values are comparable only within a method and matching polynomial specification. Missing alternatives remain visible as missing; continuation panels retain identical construction.}{G0.5 v3 modified manufacturing records with fictional company links; authors' calculations.}
\end{figure}
\begin{figure}[H]
\centering
\caption{LP: estimation sensitivity (continued)}
\GOneLabel{fig:deep-lp-sensitivity-p2}
\includegraphics[width=\linewidth]{figures/deep_lp_sensitivity_p2.pdf}
\appnotes{The primary cubic specification, two alternative starts and quadratic proxy use the same own-sample eligibility and numerical translations. Objective values are comparable only within a method and matching polynomial specification. Missing alternatives remain visible as missing; continuation panels retain identical construction.}{G0.5 v3 modified manufacturing records with fictional company links; authors' calculations.}
\end{figure}
Alternative starts change the labor coefficient by at most 0.00e+00; alternative starts change the capital coefficient by at most 1.85e-04; the quadratic proxy change the labor coefficient by at most 3.12e-02; the quadratic proxy change the capital coefficient by at most 2.06e-02. These comparisons keep each sector's records fixed. We do not use an economically convenient sign, a narrow residual distribution or a smaller criterion under a different model to select a replacement primary estimate.


\GOneSubsection{Olley--Pakes}
\GOneLabel{sec:op}
\GOneSubsubsection{Investment proxy}
\GOneLabel{sec:op-economics}
Olley and Pakes use investment to learn about unobserved production conditions and account for companies' decisions to continue operating \citep{op}. At a given capital stock, a company expecting more favorable conditions may invest more. If investment strictly increases with those conditions, after accounting for the other determinants of investment, we can work backward from investment to the unobserved component in the same way as with materials. The additional issue is \emph{selection}: the companies we still observe may differ systematically from those that disappear. For example, if less productive companies are more likely to close, the relationship between past and current productivity among continuing companies will not describe every company that could have continued.

\GOneSubsubsection{Variables and timing}
\GOneLabel{sec:op-data-mapping}
\GOneDetail{Source variables}
\GOneLabel{sec:op-source-variables}
We check the fields in Table~\GOneRef{tab:deep-op-mapping} against the requirements for OP estimation. The table links each variable in the equations to its source field and states the condition a record must meet before we take its logarithm.
\GOneInput{tables/deep_op_mapping.tex}
The outcome measures reported value added at reported prices. We measure labor by the number of workers, without information about their hours or skills. For capital, we assume that the preceding year's reported stock describes the assets available at the start of the production year. This assumption extends the calculation to the available years without establishing that their source definitions are equivalent. A stock also does not directly measure how intensively a company uses its assets. These choices affect both the estimated input contributions and the remainder we compare.

\GOneDetail{Year coverage}
\GOneLabel{sec:op-year-coverage}
For production in year $t$, we obtain beginning capital from the stock reported in calendar year $t-1$, treating that amount as its year-end stock. Materials or investment comes from year $t$ itself. The dynamic stage also needs the preceding production observation; its beginning capital consequently draws on year $t-2$. We preserve every calendar gap instead of treating the next observed year as the next calendar year. Table~\GOneRef{tab:deep-op-input-states} reports each input's availability before we require them together.
\GOneInput{tables/deep_op_input_states.tex}
The positive observations for one input need not be the same observations as those for another. We retain a record only when it supplies all the quantities that its estimation step requires. A missing response does not show zero production, and we cannot take the log of a recorded zero. Beginning capital also requires the immediately preceding stock. The year-end and investment assumptions admit positive reported amounts without changing the documentary flags; those flags continue to show where the source evidence is narrower than the assumption.

\GOneDetail{Transformations and sample restrictions}
\GOneLabel{sec:op-transformations}
We check positivity before taking logs. We use the available production years and sectors, match each record to the immediately preceding calendar year, and apply the minimum sample sizes to the pooled sector sample. For the polynomial calculations, we then subtract the relevant sample median from each logged variable; this changes its numerical reference point while keeping the eligible records fixed. We estimate one coefficient vector and control relationship per sector over the pooled years, without adding year effects. Table~\GOneRef{tab:deep-op-sample} traces the sample by sector and year.
\GOneInput{tables/deep_op_sample.tex}
The estimation steps use these records in different ways. An earlier observation can supply the preceding-year inputs for a later calendar pair. The estimators with a separate first stage can also use eligible records outside the accepted pairs in that first-stage calculation. Once a sector has estimated coefficients, we calculate a residual for each record with the required current inputs, including one without an eligible predecessor. Calculating that residual does not add a new pair or re-estimate the coefficients. An observation outside the dynamic-stage sample can therefore contribute to the first stage, receive a residual, or do both.


\GOneSubsubsection{Risk sample}
\GOneLabel{sec:op-data}
We use positive reported capital additions, assuming they measure investment as OP's investment measure. We measure beginning-of-year capital from the preceding December stock, following the LP timing rule. These two variables serve different purposes: the stock measures capital already available, while additions measure investment. Subtracting consecutive stocks would combine additions with disposals, depreciation, valuation and reporting changes. We therefore retain the reported additions measure under that assumption rather than reconstruct investment from stock differences.

We create a risk interval for every fictional company assignment at the start of each adjacent pair of supplied calendar years. If the assignment appears in the next year, we record zero for disappearance; if it has no match, we record one. We create no interval across the absent 2016 file and no interval following the terminal 2022 observation. These outcomes describe the assignment, not the survival of an actual establishment. To estimate Equation~\GOneEqref{eq:op-selection}, an interval also needs positive previous-year investment and capital. Counting every fictional ID would therefore overstate the observations available for the selection regression.

\GOneSubsubsection{Selection equation}
\GOneLabel{sec:op-stages}
To account for selective continuation, the method uses both previous productivity and the probability of disappearance when predicting current productivity among continuing companies. Let $i_{it}$ denote the logarithm of positive investment, $S_{it}=1$ identify a company that continues over the interval ending in year $t$, and $q_{it}$ denote its modeled probability of disappearance. The two relationships are
\begin{equation}
\begin{aligned}
\omega_{it}&=h(k_{it},i_{it}),\\
E[\omega_{it}\mid\omega_{i,t-1},S_{it}=1,\mathcal I_{i,t-1}]&=g_S(\omega_{i,t-1},q_{it}),
\end{aligned}
\GOneLabel{eq:op-economics}
\end{equation}
We calculate $q$ from the relevant previous-year information. We use the probability of disappearance consistently throughout the calculation; using a continuation probability would require transforming the approximation accordingly. The application also requires a credible interpretation of disappearance. Absence from a file alone does not establish that a company closed. Here the links are fictional, so their continuation describes the generated assignment rather than actual company survival.

The OP procedure first regresses logged value added on logged workers, an intercept and the capital--investment polynomial $B_3(k,i)$. It retains $\widehat\beta_L$ and constructs $\widehat\phi$ as in Equation~\GOneEqref{eq:lp-phi}. It also estimates a separate disappearance equation using all eligible company intervals, including those without a continuing production record. We specify a logistic regression, which maps the previous-year capital and investment polynomial into a probability between zero and one:
\begin{equation}
\widehat q_{it}=\Lambda\!\left(\widehat\pi_0+B_3(k_{i,t-1},i_{i,t-1})'\widehat\pi\right),
\qquad \Lambda(v)=\frac{\exp(v)}{1+\exp(v)}.
\GOneLabel{eq:op-selection}
\end{equation}
The outcome equals one when the assignment disappears and zero when it continues. The coefficients $\pi$ describe how its probability varies with previous-year capital and investment. We define this risk sample before applying the current-year production requirements: otherwise we could delete the very disappearances needed to estimate their probability. An application to actual company survival would require evidence establishing which companies could continue and which closed. Our risk sample instead follows the fictional assignment rule described above.

For each trial capital coefficient $b_K$, we calculate previous-year productivity as $w=\widehat\phi_{t-1}-b_Kk_{t-1}$. We then predict current productivity using both $w$ and the estimated disappearance probability $q$, including their powers and products up to total degree three:
\begin{equation}
H^S_{it}(b_K)=(1,w,w^2,w^3,q,q^2,q^3,qw,q^2w,qw^2).
\GOneLabel{eq:op-transition}
\end{equation}
We regress the current productivity estimate on these terms, then choose capital's coefficient by minimizing the counterpart of Equation~\GOneEqref{eq:lp-objective} among eligible continuing records. This requires a predicted disappearance probability for every production pair. The risk sample must contain both disappearances and continuations; the logistic regression must converge; and the regressors used in the selection and productivity equations must contain enough independent variation to estimate their coefficients. These last requirements are the full-rank conditions. If the selection equation fails, we do not proceed by dropping the selection correction. We retain failed or unsupported cells separately from the estimates that meet these requirements.

\GOneSubsubsection{Sample support}
\GOneLabel{sec:op-support}
We examine every observed sector in Table~\GOneRef{tab:deep-op-all-support}. This separates the availability of production inputs from the information needed to model continuation.
\GOneInput{tables/deep_op_all_support.tex}
The sectors with accepted OP estimates are 10. We report their coefficients and uncertainty below. A positive reported investment amount alone is not sufficient: the pooled production sample, continuation sample and numerical requirements must also be met.


We compare the usable investment and continuation records with the minimum sample sizes specified before estimation in Figure~\GOneRef{fig:deep-op-support}. Table~\GOneRef{tab:support} gives the corresponding production and risk counts across observed sectors. The selection equation requires 100 risk intervals, including 10 disappearances and 10 continuations. The production equations must separately meet their first-stage and consecutive-year sample requirements.
\begin{figure}[H]
\centering
\caption{Investment and continuation support}
\GOneLabel{fig:deep-op-support}
\includegraphics[width=\linewidth]{figures/deep_op_support.pdf}
\appnotes{The OP risk sample records whether eligible fictional assignments continue into the immediately following supplied calendar year, including disappearances. It is formed separately from the continuing production sample. No interval crosses the missing source year or follows the final year. Sample thresholds do not replace rank or convergence checks.}{G0.5 v3 modified manufacturing records with fictional company links; authors' calculations.}
\end{figure}
\GOneSubsubsection{Worked aggregate calculation}
\GOneLabel{sec:op-worked}
We illustrate the calculation using mean logged inputs within one sector. Multiplying each mean log input by its fitted coefficient and subtracting both contributions from mean logged value added gives the mean individual residual exactly, because the same slopes apply to every record in that sector. Table~\GOneRef{tab:deep-op-worked} follows Equations~\GOneEqref{eq:native-definition}--\GOneEqref{eq:relative-index} through centering and exponentiation, using full precision before rounding.
\GOneInput{tables/deep_op_worked.tex}
In 2011, mean logged value added is 14.4150. The fitted labor and capital contributions are 3.6738 and 1.2793, leaving 9.4619. Subtracting the pooled reference mean of 10.3468 gives a centered mean of -0.8849. Its exponential is 0.4128, while the arithmetic mean of individual index values is 1.6564. The latter gives greater weight to large individual residuals after exponentiation. Both annual examples use the same pooled equation; neither estimates a separate annual technology.


\GOneSubsubsection{Coefficient estimates}
\GOneLabel{sec:op-coefficients}
We record the primary outcome for every sector and both input-sample definitions in Table~\GOneRef{tab:deep-op-attempts}. This keeps the absence of a fit separate from a coefficient estimated near zero.
\GOneInput{tables/deep_op_attempts.tex}
Only estimated cells supply coefficients and derived residuals. A count requirement and a numerical failure have different meanings: the former prevents the declared calculation from being attempted, while the latter prevents an attempted calculation from supplying an accepted estimate.
We report the estimated coefficients beginning with Table~\GOneRef{tab:deep-op-regression}. The own and common samples have separate columns and fits. To describe uncertainty, we resample whole fictional company histories and repeat every estimation step. Standard errors measure the spread of these repeated estimates; percentile intervals give their admitted lower and upper quantiles.
\GOneInput{tables/deep_op_regression.tex}
\GOneInput{tables/deep_op_regression-common.tex}
The uncertainty in the coefficient sum depends on how labor and capital estimates move together across resamples. Adding their separate standard errors would ignore that covariance. Fitted year coverage is pooled, and a separate first stage can use more records than the dynamic sample. These counts describe distinct roles in the same estimation.

Table~\GOneRef{tab:deep-op-numerics} reports the rank and optimization diagnostics for the successful primary fits. We compare numerical rank with the dimension of the relevant matrix because returning a coefficient is not enough to show that the stated calculation had the required independent information.
\GOneInput{tables/deep_op_numerics.tex}
The diagnostics describe the numerical problem at the reported solution. They do not establish valid economic timing or a unique global optimum; the alternative-start comparison examines a separate aspect of that uncertainty.

Figure~\GOneRef{fig:deep-op-coefficients} displays the same own-sample coefficients and full admitted uncertainty intervals. Separate input panels preserve the scale of the labor and capital estimates.
\begin{figure}[H]
\centering
\caption{OP: coefficient estimates}
\GOneLabel{fig:deep-op-coefficients}
\includegraphics[width=\linewidth]{figures/deep_op_coefficients.pdf}
\appnotes{We pool supported years within each sector under constant coefficients. Marks show own-sample primary estimates and intervals admitted under the bootstrap minimum-success rule. Intervals retain their full endpoints and condition on the fictional histories and maintained measurements.}{G0.5 v3 modified manufacturing records with fictional company links; authors' calculations.}
\end{figure}
In sector 10, the labor coefficient is 0.8942; the capital coefficient is 0.0908. We retain the signs and magnitudes returned by the primary specification. Their interpretation depends on beginning-capital timing, reported monetary scales, the proxy assumptions and the estimation criterion. The intervals describe additional variation when the same calculation is repeated on resampled histories.
\GOneSubsubsection{Bootstrap coefficient distributions}
\GOneLabel{sec:op-bootstrap}
Table~\GOneRef{tab:deep-op-bootstrap} reports the spread of the repeated coefficient estimates. The standard error is their sample standard deviation, while the median and percentiles reveal features that a single standard error can conceal. The minimum and maximum make finite extremes visible.
\GOneInput{tables/deep_op_bootstrap.tex}
The largest labor-coefficient standard deviation occurs in sector 10, at 0.0613. Its bootstrap median is 0.8989, with P5 and P95 at 0.7828 and 0.9983; the full range is 0.7401 to 1.0251. Comparing the middle of this distribution with its extremes helps distinguish broad dispersion from a small number of distant estimates. We retain every finite resampled estimate and keep the original-sample coefficient as the primary result. A limited number of resamples determines the extreme percentiles, so their locations should not be interpreted as precise tail probabilities. The exercise also leaves uncertainty about actual identities, prices and economic restrictions outside the resampling calculation.
\GOneSubsubsection{Residual distributions}
\GOneLabel{sec:op-distribution}
We subtract the fitted labor and capital contributions from logged value added, then subtract each sector's pooled dynamic residual mean. Figure~\GOneRef{fig:deep-op-histogram} compares the uncentered and centered distributions on exactly the same fitted records.
\begin{figure}[H]
\centering
\caption{OP: pooled residual distributions}
\GOneLabel{fig:deep-op-histogram}
\includegraphics[width=\linewidth]{figures/deep_op_histogram.pdf}
\appnotes{Both panels contain the same own-sample dynamic observations pooled across all fitted years. Uncentered residuals use original logs; centering subtracts one mean per sector-method-sample. Bins span 0.5 log points and retain all finite tails. Labels show the modal and first and last occupied bins with source-derived counts and shares.}{G0.5 v3 modified manufacturing records with fictional company links; authors' calculations.}
\end{figure}
The centered distribution has a standard deviation of 1.363 log points. Centering removes differences in sector means while preserving all within-sector gaps and rankings. The pooled shape can change because it previously combined those mean differences with variation within sectors. We retain all finite observations, including the tails.
\GOneSubsubsection{Prediction samples and tails}
\GOneLabel{sec:op-prediction}
The histogram above describes the dynamic estimation sample. We can calculate residuals more broadly whenever a record supplies the required current inputs in a sector with estimated slopes. Table~\GOneRef{tab:deep-op-scopes} distinguishes these roles for every source year and for the pooled records.
\GOneInput{tables/deep_op_scopes.tex}
There are 683 predicted company-year observations: 128 belong to the pooled dynamic fit and 555 lie outside it. Only the dynamic-fit records determine the reference mean. The annual fit means and the broader prediction means can therefore differ from zero, even though the pooled fit mean is zero by construction. A company appearing repeatedly contributes several observations to these counts.

Table~\GOneRef{tab:deep-op-quantiles} holds the fitted records fixed and compares the uncentered logs, centered logs and positive index. We include the minimum and maximum so that the transformed tails remain visible.
\GOneInput{tables/deep_op_quantiles.tex}
Centering preserves every within-sector gap. Exponentiation preserves within-sector ordering but enlarges distances at high values, as Equation~\GOneEqref{eq:relative-index} shows. A large index maximum therefore needs to be read through this transformation, rather than treated directly as physical efficiency.

Table~\GOneRef{tab:deep-op-tails} identifies distant observations using two distinct rules. The interquartile range (IQR) is the distance from P25 to P75. Fences placed 1.5 IQR beyond those quartiles compare a value with the spread of its own distribution. Fixed absolute-log and index cutoffs instead apply the same numerical threshold in every sector. A record can exceed more than one cutoff.
\GOneInput{tables/deep_op_tails.tex}
In the pooled dynamic sample, 13 records have an absolute centered log residual above two. These flags describe the residuals produced by the fitted equation. Prices, reported amounts, the estimated input contributions and production conditions can all affect an extreme residual, so the flag alone identifies neither an error nor its cause. We retain every flagged observation.

\GOneSubsubsection{Annual distributions}
\GOneLabel{sec:op-annual}
Table~\GOneRef{tab:deep-op-annual} compares all eligible predictions by sector and year. We keep each sector's fitted slopes and reference mean constant over time. This preserves annual differences instead of subtracting each year's average separately.
\GOneInput{tables/deep_op_annual.tex}
Changes in annual means and quartiles combine differences in reported outcomes and inputs with differences in sample composition. A year with more high-residual observations can have a higher mean without every continuing assignment improving. We examine this distinction next by matching observations within adjacent years.
Figure~\GOneRef{fig:deep-op-year-sector} and its continuation panels plot the annual median and quartiles for each sector. These intervals measure dispersion across observations, not uncertainty in an estimated coefficient.
\begin{figure}[H]
\centering
\caption{OP annual residual quartiles}
\GOneLabel{fig:deep-op-year-sector}
\includegraphics[width=\linewidth]{figures/deep_op_year_sector.pdf}
\appnotes{Each sector panel includes the available annual predictions, using one pooled coefficient vector and reference mean. Dots mark medians and intervals run from P25 to P75; all available values are numerically labelled. An empty group is not a zero. The legend identifies any missing-file region. Continuation panels use the same construction.}{G0.5 v3 modified manufacturing records with fictional company links; authors' calculations.}
\end{figure}

The sector panels separate variation within each industry from changes in which industries dominate a pooled annual distribution. They describe the same fixed sector equations over time.
\GOneSubsubsection{Adjacent-year changes}
\GOneLabel{sec:op-paired}
We hold membership fixed within each comparison by requiring a prediction for the same fictional assignment and sector in adjacent calendar years. Table~\GOneRef{tab:deep-op-pair-coverage} shows how this restriction narrows the endpoint samples. We retain intervals with no matched predictions so that missing capital history remains visible.
\GOneInput{tables/deep_op_pair_coverage.tex}
The pooled comparison contains 128 record-pair intervals. A company can enter several successive intervals, so this is not a count of independent companies. Restricting to matched predictions removes changes in membership within each interval, although membership can still differ between intervals.

For each pair, we subtract the earlier centered residual from the later one. Let $\Delta x_{it}=x_{it}-x_{i,t-1}$ for any logged variable $x$. Holding the sector slopes fixed gives
\begin{equation}
\Delta\widetilde r_{it}=\Delta y_{it}-\widehat\beta_{L,s}\Delta l_{it}-\widehat\beta_{K,s}\Delta k_{it}.
\GOneLabel{eq:op-paired-change}
\end{equation}
The pooled reference mean cancels because we subtract the same constant in both years. Table~\GOneRef{tab:deep-op-paired} reports the resulting changes for each interval and sector, then pools the record-pair intervals under the same rule.
\GOneInput{tables/deep_op_paired.tex}
The pooled mean log change is 0.096. The correlation of earlier and later ranks is 0.682, measuring persistence in ordering. The mean individual index ratio is 3.573; large ratios can strongly influence this arithmetic mean, which need not equal the exponential of the mean log change. These comparisons use fictional links and reported-price equations; they do not establish real-company productivity changes.
Figure~\GOneRef{fig:deep-op-paired-change} and its continuation panels plot these adjacent-year changes under the same matched-prediction rule.
\begin{figure}[H]
\centering
\caption{OP adjacent-year residual changes}
\GOneLabel{fig:deep-op-paired-change}
\includegraphics[width=\linewidth]{figures/deep_op_paired_change.pdf}
\appnotes{Each comparison uses same-sector fictional assignments with eligible predictions in consecutive calendar years. Sector coefficients and pooled reference means remain fixed. Histograms retain all finite changes. Numeric labels identify the modal and first and last occupied bins. Pooled panels combine record-pair intervals, allowing repeated company contributions; interval panels retain their own denominators.}{G0.5 v3 modified manufacturing records with fictional company links; authors' calculations.}
\end{figure}
\begin{figure}[H]
\centering
\caption{OP adjacent-year residual changes (continued)}
\GOneLabel{fig:deep-op-paired-change-p2}
\includegraphics[width=\linewidth]{figures/deep_op_paired_change_p2.pdf}
\appnotes{Each comparison uses same-sector fictional assignments with eligible predictions in consecutive calendar years. Sector coefficients and pooled reference means remain fixed. Histograms retain all finite changes. Numeric labels identify the modal and first and last occupied bins. Pooled panels combine record-pair intervals, allowing repeated company contributions; interval panels retain their own denominators.}{G0.5 v3 modified manufacturing records with fictional company links; authors' calculations.}
\end{figure}

A pooled average can conceal increases and decreases across matched assignments. The interval panels show how their distributions vary over time under the fixed sector equations.
\GOneSubsubsection{Specification sensitivity}
\GOneLabel{sec:op-sensitivity}
Figure~\GOneRef{fig:deep-op-sensitivity} and its continuation panels compare the prespecified starting values and proxy-polynomial alternatives. Changing the start leaves the model fixed; changing the polynomial alters the fitted control relationship. A numerical stopping rule need not identify the same solution from every starting value.
\begin{figure}[H]
\centering
\caption{OP: estimation sensitivity}
\GOneLabel{fig:deep-op-sensitivity}
\includegraphics[width=\linewidth]{figures/deep_op_sensitivity.pdf}
\appnotes{The primary cubic specification, two alternative starts and quadratic proxy use the same own-sample eligibility and numerical translations. Objective values are comparable only within a method and matching polynomial specification. Missing alternatives remain visible as missing; continuation panels retain identical construction.}{G0.5 v3 modified manufacturing records with fictional company links; authors' calculations.}
\end{figure}
Alternative starts change the labor coefficient by at most 0.00e+00; alternative starts change the capital coefficient by at most 9.15e-05; the quadratic proxy change the labor coefficient by at most 2.42e-02; the quadratic proxy change the capital coefficient by at most 1.82e-01. These comparisons keep each sector's records fixed. We do not use an economically convenient sign, a narrow residual distribution or a smaller criterion under a different model to select a replacement primary estimate.


Investment changes both the proxy information and the continuation adjustment. A separate issue concerns whether labor's contribution can be distinguished from the information carried by materials, even with an adequate sample. Ackerberg--Caves--Frazer addresses that issue.

\GOneSubsection{Ackerberg--Caves--Frazer}
\GOneLabel{sec:acf}
\GOneSubsubsection{Labor identification}
\GOneLabel{sec:acf-economics}
Ackerberg, Caves and Frazer examine whether LP's first-stage regression can separate labor's contribution from the capital--materials function \citep{acf}. Suppose a company chooses labor and materials using the same production information. Employment may then be predictable from the very capital and materials variables used to account for unobserved conditions. The regression would have no independent variation with which to distinguish labor's effect from that function. A finite polynomial can still return a labor coefficient, but the coefficient would depend on the approximation rather than establish the separate economic contribution of labor. This is the labor-identification problem.

We follow ACF by leaving labor and capital combined with the predictable part of production in the first stage. We then estimate both input coefficients together in the second stage, using restrictions on the new productivity component that remains after accounting for the previous year. This shifts the work of separating the two input contributions to the second stage. The timing assumptions and the information available in the selected variables still determine whether that separation is credible.

\GOneSubsubsection{Variables and timing}
\GOneLabel{sec:acf-data}
We use the same positive reported nominal value added, workers, preceding-December capital and materials measures as for LP. We apply the same eligibility rules and require consecutive calendar years for the second stage. ACF's first-stage regression nevertheless differs because its polynomial includes labor and labor's products with capital and materials. We estimate the two input contributions jointly on all eligible consecutive-year records pooled within a sector. We also apply these coefficients to observations with eligible current inputs outside the dynamic stage. Those observations may already have contributed to the first-stage regression; predicting their residuals does not fit another annual equation.

Using the same observed measures lets us trace the residual difference between ACF and LP to their coefficients. Within a sector, subtracting the LP residual from the ACF residual gives the negative of the labor-coefficient difference times logged workers, plus the negative of the capital-coefficient difference times logged capital. Centering then subtracts that difference's mean in the sector's second-stage sample. We therefore examine the estimated coefficients before interpreting how ACF changes the residual distribution.

\GOneDetail{Source variables}
\GOneLabel{sec:acf-source-variables}
We check the fields in Table~\GOneRef{tab:deep-acf-mapping} against the requirements for ACF estimation. The table links each variable in the equations to its source field and states the condition a record must meet before we take its logarithm.
\GOneInput{tables/deep_acf_mapping.tex}
The outcome measures reported value added at reported prices. We measure labor by the number of workers, without information about their hours or skills. For capital, we assume that the preceding year's reported stock describes the assets available at the start of the production year. This assumption extends the calculation to the available years without establishing that their source definitions are equivalent. A stock also does not directly measure how intensively a company uses its assets. These choices affect both the estimated input contributions and the remainder we compare.

\GOneDetail{Year coverage}
\GOneLabel{sec:acf-year-coverage}
For production in year $t$, we obtain beginning capital from the stock reported in calendar year $t-1$, treating that amount as its year-end stock. Materials or investment comes from year $t$ itself. The dynamic stage also needs the preceding production observation; its beginning capital consequently draws on year $t-2$. We preserve every calendar gap instead of treating the next observed year as the next calendar year. Table~\GOneRef{tab:deep-acf-input-states} reports each input's availability before we require them together.
\GOneInput{tables/deep_acf_input_states.tex}
The positive observations for one input need not be the same observations as those for another. We retain a record only when it supplies all the quantities that its estimation step requires. A missing response does not show zero production, and we cannot take the log of a recorded zero. Beginning capital also requires the immediately preceding stock. The year-end and investment assumptions admit positive reported amounts without changing the documentary flags; those flags continue to show where the source evidence is narrower than the assumption.

\GOneDetail{Transformations and sample restrictions}
\GOneLabel{sec:acf-transformations}
We check positivity before taking logs. We use the available production years and sectors, match each record to the immediately preceding calendar year, and apply the minimum sample sizes to the pooled sector sample. For the polynomial calculations, we then subtract the relevant sample median from each logged variable; this changes its numerical reference point while keeping the eligible records fixed. We estimate one coefficient vector and control relationship per sector over the pooled years, without adding year effects. Table~\GOneRef{tab:deep-acf-sample} traces the sample by sector and year.
\GOneInput{tables/deep_acf_sample.tex}
The estimation steps use these records in different ways. An earlier observation can supply the preceding-year inputs for a later calendar pair. The estimators with a separate first stage can also use eligible records outside the accepted pairs in that first-stage calculation. Once a sector has estimated coefficients, we calculate a residual for each record with the required current inputs, including one without an eligible predecessor. Calculating that residual does not add a new pair or re-estimate the coefficients. An observation outside the dynamic-stage sample can therefore contribute to the first stage, receive a residual, or do both.


\GOneSubsubsection{Estimation procedure}
\GOneLabel{sec:acf-stages}
We first predict logged value added from logged workers, capital and materials without assigning a separate economic coefficient to labor:
\begin{equation}
y_{it}=\Phi(l_{it},k_{it},m_{it})+\varepsilon_{it},\qquad
\widehat\Phi_{it}=\widehat E[y_{it}\mid l_{it},k_{it},m_{it}],
\GOneLabel{eq:acf-first}
\end{equation}
We use an intercept and every power and product of $(l,k,m)$ with total degree up to three. The predicted value $\widehat\Phi$ includes the contributions of labor and capital, so it is not yet an estimate of productivity alone. For each trial pair of coefficients $b=(b_L,b_K)'$, we subtract the implied input contributions. We then subtract the part of this provisional productivity estimate that we can predict from its previous-year value:
\begin{equation}
\widehat\omega_{it}(b)=\widehat\Phi_{it}-b_Ll_{it}-b_Kk_{it},\qquad
\widehat\xi_{it}(b)=\widehat\omega_{it}(b)-\widehat g_b\!\left(\widehat\omega_{i,t-1}(b)\right).
\GOneLabel{eq:acf-innovation}
\end{equation}
To obtain $\widehat g_b$, we regress current provisional productivity on an intercept and the first three powers of its previous-year value. We refit this regression for every trial pair of input coefficients. The remaining component, $\widehat\xi(b)$, estimates the current-year productivity surprise under that trial pair.

We choose the coefficient pair by asking whether the remaining surprise is related to information that, under our timing assumptions, should not predict it. The two variables we use are previous-year labor and beginning-of-year capital. They serve as \emph{instruments}: their assumed lack of association with the surprise supplies restrictions for estimating the coefficients. For each instrument, we form a \emph{moment} by multiplying it by $\widehat\xi(b)$ and averaging across the estimation sample. The assumed population moments are zero. With $Z=[l_{t-1},k_t]$ and $N$ consecutive-year observations, we minimize the following weighted measure of how far the sample moments are from zero:
\begin{equation}
Q_{\mathrm{ACF}}(b)=\left[Z'\widehat\xi(b)\right]'
\frac{(Z'Z)^{-1}}{N}\left[Z'\widehat\xi(b)\right].
\GOneLabel{eq:acf}
\end{equation}
The zero-product restrictions are also described as \emph{orthogonality} between the instruments and the productivity innovation. We check that the instruments contain independent information and that changes in each input coefficient produce distinguishable changes in the moments near the estimated solution. These are the instrument-rank and moment-derivative-rank checks. They assess whether we can solve the numerical problem, not whether company decisions satisfy the assumed timing or materials-demand rule. With two moments for two input coefficients, a value near zero for the criterion does not provide an additional, independent test of those assumptions. Fictional company histories also cannot establish their validity for actual companies.

\GOneSubsubsection{Worked aggregate calculation}
\GOneLabel{sec:acf-worked}
We illustrate the calculation using mean logged inputs within one sector. Multiplying each mean log input by its fitted coefficient and subtracting both contributions from mean logged value added gives the mean individual residual exactly, because the same slopes apply to every record in that sector. Table~\GOneRef{tab:deep-acf-worked} follows Equations~\GOneEqref{eq:native-definition}--\GOneEqref{eq:relative-index} through centering and exponentiation, using full precision before rounding.
\GOneInput{tables/deep_acf_worked.tex}
In 2011, mean logged value added is 13.7219. The fitted labor and capital contributions are 4.3106 and 1.9360, leaving 7.4753. Subtracting the pooled reference mean of 8.1934 gives a centered mean of -0.7180. Its exponential is 0.4877, while the arithmetic mean of individual index values is 1.4169. The latter gives greater weight to large individual residuals after exponentiation. Both annual examples use the same pooled equation; neither estimates a separate annual technology.


\GOneSubsubsection{Coefficient estimates}
\GOneLabel{sec:acf-coefficients}
We record the primary outcome for every sector and both input-sample definitions in Table~\GOneRef{tab:deep-acf-attempts}. This keeps the absence of a fit separate from a coefficient estimated near zero.
\GOneInput{tables/deep_acf_attempts.tex}
Only estimated cells supply coefficients and derived residuals. A count requirement and a numerical failure have different meanings: the former prevents the declared calculation from being attempted, while the latter prevents an attempted calculation from supplying an accepted estimate.
We report the estimated coefficients beginning with Table~\GOneRef{tab:deep-acf-regression}. The own and common samples have separate columns and fits. To describe uncertainty, we resample whole fictional company histories and repeat every estimation step. Standard errors measure the spread of these repeated estimates; percentile intervals give their admitted lower and upper quantiles.
\GOneInput{tables/deep_acf_regression.tex}
\GOneInput{tables/deep_acf_regression-2.tex}
\GOneInput{tables/deep_acf_regression-3.tex}
\GOneInput{tables/deep_acf_regression-common.tex}
The uncertainty in the coefficient sum depends on how labor and capital estimates move together across resamples. Adding their separate standard errors would ignore that covariance. Fitted year coverage is pooled, and a separate first stage can use more records than the dynamic sample. These counts describe distinct roles in the same estimation.

Table~\GOneRef{tab:deep-acf-numerics} reports the rank and optimization diagnostics for the successful primary fits. We compare numerical rank with the dimension of the relevant matrix because returning a coefficient is not enough to show that the stated calculation had the required independent information.
\GOneInput{tables/deep_acf_numerics.tex}
The diagnostics describe the numerical problem at the reported solution. They do not establish valid economic timing or a unique global optimum; the alternative-start comparison examines a separate aspect of that uncertainty.

Figure~\GOneRef{fig:deep-acf-coefficients} displays the same own-sample coefficients and full admitted uncertainty intervals. Separate input panels preserve the scale of the labor and capital estimates.
\begin{figure}[H]
\centering
\caption{ACF: coefficient estimates}
\GOneLabel{fig:deep-acf-coefficients}
\includegraphics[width=\linewidth]{figures/deep_acf_coefficients.pdf}
\appnotes{We pool supported years within each sector under constant coefficients. Marks show own-sample primary estimates and intervals admitted under the bootstrap minimum-success rule. Intervals retain their full endpoints and condition on the fictional histories and maintained measurements.}{G0.5 v3 modified manufacturing records with fictional company links; authors' calculations.}
\end{figure}
Across the estimated sectors, labor ranges from 0.9675 to 1.4197; capital ranges from 0.0113 to 0.1735. We retain the signs and magnitudes returned by the primary specification. Their interpretation depends on beginning-capital timing, reported monetary scales, the proxy assumptions and the estimation criterion. The intervals describe additional variation when the same calculation is repeated on resampled histories.
\GOneSubsubsection{Bootstrap coefficient distributions}
\GOneLabel{sec:acf-bootstrap}
Table~\GOneRef{tab:deep-acf-bootstrap} reports the spread of the repeated coefficient estimates. The standard error is their sample standard deviation, while the median and percentiles reveal features that a single standard error can conceal. The minimum and maximum make finite extremes visible.
\GOneInput{tables/deep_acf_bootstrap.tex}
The largest labor-coefficient standard deviation occurs in sector 15, at 3.8544. Its bootstrap median is 1.1351, with P5 and P95 at 1.0315 and 10.5967; the full range is 1.0114 to 26.9816. Comparing the middle of this distribution with its extremes helps distinguish broad dispersion from a small number of distant estimates. We retain every finite resampled estimate and keep the original-sample coefficient as the primary result. A limited number of resamples determines the extreme percentiles, so their locations should not be interpreted as precise tail probabilities. The exercise also leaves uncertainty about actual identities, prices and economic restrictions outside the resampling calculation.
\GOneSubsubsection{Residual distributions}
\GOneLabel{sec:acf-distribution}
We subtract the fitted labor and capital contributions from logged value added, then subtract each sector's pooled dynamic residual mean. Figure~\GOneRef{fig:deep-acf-histogram} compares the uncentered and centered distributions on exactly the same fitted records.
\begin{figure}[H]
\centering
\caption{ACF: pooled residual distributions}
\GOneLabel{fig:deep-acf-histogram}
\includegraphics[width=\linewidth]{figures/deep_acf_histogram.pdf}
\appnotes{Both panels contain the same own-sample dynamic observations pooled across all fitted years. Uncentered residuals use original logs; centering subtracts one mean per sector-method-sample. Bins span 0.5 log points and retain all finite tails. Labels show the modal and first and last occupied bins with source-derived counts and shares.}{G0.5 v3 modified manufacturing records with fictional company links; authors' calculations.}
\end{figure}
The centered distribution has a standard deviation of 1.130 log points. Centering removes differences in sector means while preserving all within-sector gaps and rankings. The pooled shape can change because it previously combined those mean differences with variation within sectors. We retain all finite observations, including the tails.
\GOneSubsubsection{Prediction samples and tails}
\GOneLabel{sec:acf-prediction}
The histogram above describes the dynamic estimation sample. We can calculate residuals more broadly whenever a record supplies the required current inputs in a sector with estimated slopes. Table~\GOneRef{tab:deep-acf-scopes} distinguishes these roles for every source year and for the pooled records.
\GOneInput{tables/deep_acf_scopes.tex}
There are 10,315 predicted company-year observations: 5,041 belong to the pooled dynamic fit and 5,274 lie outside it. Only the dynamic-fit records determine the reference mean. The annual fit means and the broader prediction means can therefore differ from zero, even though the pooled fit mean is zero by construction. A company appearing repeatedly contributes several observations to these counts.

Table~\GOneRef{tab:deep-acf-quantiles} holds the fitted records fixed and compares the uncentered logs, centered logs and positive index. We include the minimum and maximum so that the transformed tails remain visible.
\GOneInput{tables/deep_acf_quantiles.tex}
Centering preserves every within-sector gap. Exponentiation preserves within-sector ordering but enlarges distances at high values, as Equation~\GOneEqref{eq:relative-index} shows. A large index maximum therefore needs to be read through this transformation, rather than treated directly as physical efficiency.

Table~\GOneRef{tab:deep-acf-tails} identifies distant observations using two distinct rules. The interquartile range (IQR) is the distance from P25 to P75. Fences placed 1.5 IQR beyond those quartiles compare a value with the spread of its own distribution. Fixed absolute-log and index cutoffs instead apply the same numerical threshold in every sector. A record can exceed more than one cutoff.
\GOneInput{tables/deep_acf_tails.tex}
In the pooled dynamic sample, 377 records have an absolute centered log residual above two. These flags describe the residuals produced by the fitted equation. Prices, reported amounts, the estimated input contributions and production conditions can all affect an extreme residual, so the flag alone identifies neither an error nor its cause. We retain every flagged observation.

\GOneSubsubsection{Annual distributions}
\GOneLabel{sec:acf-annual}
Table~\GOneRef{tab:deep-acf-annual} compares all eligible predictions by sector and year. We keep each sector's fitted slopes and reference mean constant over time. This preserves annual differences instead of subtracting each year's average separately.
\GOneInput{tables/deep_acf_annual.tex}
Changes in annual means and quartiles combine differences in reported outcomes and inputs with differences in sample composition. A year with more high-residual observations can have a higher mean without every continuing assignment improving. We examine this distinction next by matching observations within adjacent years.
Figure~\GOneRef{fig:deep-acf-year-sector} and its continuation panels plot the annual median and quartiles for each sector. These intervals measure dispersion across observations, not uncertainty in an estimated coefficient.
\begin{figure}[H]
\centering
\caption{ACF annual residual quartiles}
\GOneLabel{fig:deep-acf-year-sector}
\includegraphics[width=\linewidth]{figures/deep_acf_year_sector.pdf}
\appnotes{Each sector panel includes the available annual predictions, using one pooled coefficient vector and reference mean. Dots mark medians and intervals run from P25 to P75; all available values are numerically labelled. An empty group is not a zero. The legend identifies any missing-file region. Continuation panels use the same construction.}{G0.5 v3 modified manufacturing records with fictional company links; authors' calculations.}
\end{figure}
\begin{figure}[H]
\centering
\caption{ACF annual residual quartiles (continued)}
\GOneLabel{fig:deep-acf-year-sector-p2}
\includegraphics[width=\linewidth]{figures/deep_acf_year_sector_p2.pdf}
\appnotes{Each sector panel includes the available annual predictions, using one pooled coefficient vector and reference mean. Dots mark medians and intervals run from P25 to P75; all available values are numerically labelled. An empty group is not a zero. The legend identifies any missing-file region. Continuation panels use the same construction.}{G0.5 v3 modified manufacturing records with fictional company links; authors' calculations.}
\end{figure}
\begin{figure}[H]
\centering
\caption{ACF annual residual quartiles (continued)}
\GOneLabel{fig:deep-acf-year-sector-p3}
\includegraphics[width=\linewidth]{figures/deep_acf_year_sector_p3.pdf}
\appnotes{Each sector panel includes the available annual predictions, using one pooled coefficient vector and reference mean. Dots mark medians and intervals run from P25 to P75; all available values are numerically labelled. An empty group is not a zero. The legend identifies any missing-file region. Continuation panels use the same construction.}{G0.5 v3 modified manufacturing records with fictional company links; authors' calculations.}
\end{figure}

The sector panels separate variation within each industry from changes in which industries dominate a pooled annual distribution. They describe the same fixed sector equations over time.
\GOneSubsubsection{Adjacent-year changes}
\GOneLabel{sec:acf-paired}
We hold membership fixed within each comparison by requiring a prediction for the same fictional assignment and sector in adjacent calendar years. Table~\GOneRef{tab:deep-acf-pair-coverage} shows how this restriction narrows the endpoint samples. We retain intervals with no matched predictions so that missing capital history remains visible.
\GOneInput{tables/deep_acf_pair_coverage.tex}
The pooled comparison contains 5,041 record-pair intervals. A company can enter several successive intervals, so this is not a count of independent companies. Restricting to matched predictions removes changes in membership within each interval, although membership can still differ between intervals.

For each pair, we subtract the earlier centered residual from the later one. Let $\Delta x_{it}=x_{it}-x_{i,t-1}$ for any logged variable $x$. Holding the sector slopes fixed gives
\begin{equation}
\Delta\widetilde r_{it}=\Delta y_{it}-\widehat\beta_{L,s}\Delta l_{it}-\widehat\beta_{K,s}\Delta k_{it}.
\GOneLabel{eq:acf-paired-change}
\end{equation}
The pooled reference mean cancels because we subtract the same constant in both years. Table~\GOneRef{tab:deep-acf-paired} reports the resulting changes for each interval and sector, then pools the record-pair intervals under the same rule.
\GOneInput{tables/deep_acf_paired.tex}
The pooled mean log change is 0.154. The correlation of earlier and later ranks is 0.588, measuring persistence in ordering. The mean individual index ratio is 2.090; large ratios can strongly influence this arithmetic mean, which need not equal the exponential of the mean log change. These comparisons use fictional links and reported-price equations; they do not establish real-company productivity changes.
Figure~\GOneRef{fig:deep-acf-paired-change} and its continuation panels plot these adjacent-year changes under the same matched-prediction rule.
\begin{figure}[H]
\centering
\caption{ACF adjacent-year residual changes}
\GOneLabel{fig:deep-acf-paired-change}
\includegraphics[width=\linewidth]{figures/deep_acf_paired_change.pdf}
\appnotes{Each comparison uses same-sector fictional assignments with eligible predictions in consecutive calendar years. Sector coefficients and pooled reference means remain fixed. Histograms retain all finite changes. Numeric labels identify the modal and first and last occupied bins. Pooled panels combine record-pair intervals, allowing repeated company contributions; interval panels retain their own denominators.}{G0.5 v3 modified manufacturing records with fictional company links; authors' calculations.}
\end{figure}
\begin{figure}[H]
\centering
\caption{ACF adjacent-year residual changes (continued)}
\GOneLabel{fig:deep-acf-paired-change-p2}
\includegraphics[width=\linewidth]{figures/deep_acf_paired_change_p2.pdf}
\appnotes{Each comparison uses same-sector fictional assignments with eligible predictions in consecutive calendar years. Sector coefficients and pooled reference means remain fixed. Histograms retain all finite changes. Numeric labels identify the modal and first and last occupied bins. Pooled panels combine record-pair intervals, allowing repeated company contributions; interval panels retain their own denominators.}{G0.5 v3 modified manufacturing records with fictional company links; authors' calculations.}
\end{figure}
\begin{figure}[H]
\centering
\caption{ACF adjacent-year residual changes (continued)}
\GOneLabel{fig:deep-acf-paired-change-p3}
\includegraphics[width=\linewidth]{figures/deep_acf_paired_change_p3.pdf}
\appnotes{Each comparison uses same-sector fictional assignments with eligible predictions in consecutive calendar years. Sector coefficients and pooled reference means remain fixed. Histograms retain all finite changes. Numeric labels identify the modal and first and last occupied bins. Pooled panels combine record-pair intervals, allowing repeated company contributions; interval panels retain their own denominators.}{G0.5 v3 modified manufacturing records with fictional company links; authors' calculations.}
\end{figure}
\begin{figure}[H]
\centering
\caption{ACF adjacent-year residual changes (continued)}
\GOneLabel{fig:deep-acf-paired-change-p4}
\includegraphics[width=\linewidth]{figures/deep_acf_paired_change_p4.pdf}
\appnotes{Each comparison uses same-sector fictional assignments with eligible predictions in consecutive calendar years. Sector coefficients and pooled reference means remain fixed. Histograms retain all finite changes. Numeric labels identify the modal and first and last occupied bins. Pooled panels combine record-pair intervals, allowing repeated company contributions; interval panels retain their own denominators.}{G0.5 v3 modified manufacturing records with fictional company links; authors' calculations.}
\end{figure}

A pooled average can conceal increases and decreases across matched assignments. The interval panels show how their distributions vary over time under the fixed sector equations.
\GOneSubsubsection{Specification sensitivity}
\GOneLabel{sec:acf-sensitivity}
Figure~\GOneRef{fig:deep-acf-sensitivity} and its continuation panels compare the prespecified starting values and proxy-polynomial alternatives. Changing the start leaves the model fixed; changing the polynomial alters the fitted control relationship. A numerical stopping rule need not identify the same solution from every starting value.
\begin{figure}[H]
\centering
\caption{ACF: estimation sensitivity}
\GOneLabel{fig:deep-acf-sensitivity}
\includegraphics[width=\linewidth]{figures/deep_acf_sensitivity.pdf}
\appnotes{The primary cubic specification, two alternative starts and quadratic proxy use the same own-sample eligibility and numerical translations. Objective values are comparable only within a method and matching polynomial specification. Missing alternatives remain visible as missing; continuation panels retain identical construction.}{G0.5 v3 modified manufacturing records with fictional company links; authors' calculations.}
\end{figure}
\begin{figure}[H]
\centering
\caption{ACF: estimation sensitivity (continued)}
\GOneLabel{fig:deep-acf-sensitivity-p2}
\includegraphics[width=\linewidth]{figures/deep_acf_sensitivity_p2.pdf}
\appnotes{The primary cubic specification, two alternative starts and quadratic proxy use the same own-sample eligibility and numerical translations. Objective values are comparable only within a method and matching polynomial specification. Missing alternatives remain visible as missing; continuation panels retain identical construction.}{G0.5 v3 modified manufacturing records with fictional company links; authors' calculations.}
\end{figure}
Alternative starts change the labor coefficient by at most 9.26e-01; alternative starts change the capital coefficient by at most 2.55e-02; the quadratic proxy change the labor coefficient by at most 1.23e-01; the quadratic proxy change the capital coefficient by at most 1.81e-02. These comparisons keep each sector's records fixed. We do not use an economically convenient sign, a narrow residual distribution or a smaller criterion under a different model to select a replacement primary estimate.


\GOneSubsection{Wooldridge}
\GOneLabel{sec:wrdg}
Wooldridge combines two ways of explaining the same observed output in one estimation procedure \citep{wooldridge}. In the first equation, we use current capital and materials purchases to account for production conditions known to the company. What remains is the discrepancy $\varepsilon$ that the company did not know when choosing its inputs. In the second equation, we use previous-year capital and materials to predict current conditions. That prediction cannot account for the current productivity surprise $\xi$, so the remaining error includes both $\xi$ and $\varepsilon$. The difference in these errors determines which variables can supply valid instruments for each equation.

\GOneSubsubsection{Variables and timing}
\GOneLabel{sec:wrdg-data}
We use the reported nominal value added, workers, preceding-December capital and materials measures already described for LP and ACF. Wooldridge therefore changes the restrictions used to estimate their contributions while keeping the observed measures fixed. Its two equations use current and previous-year values for the same eligible calendar pairs, pooling all supported years within each sector.

\GOneDetail{Source variables}
\GOneLabel{sec:wrdg-source-variables}
We check the fields in Table~\GOneRef{tab:deep-wrdg-mapping} against the requirements for Wooldridge estimation. The table links each variable in the equations to its source field and states the condition a record must meet before we take its logarithm.
\GOneInput{tables/deep_wrdg_mapping.tex}
The outcome measures reported value added at reported prices. We measure labor by the number of workers, without information about their hours or skills. For capital, we assume that the preceding year's reported stock describes the assets available at the start of the production year. This assumption extends the calculation to the available years without establishing that their source definitions are equivalent. A stock also does not directly measure how intensively a company uses its assets. These choices affect both the estimated input contributions and the remainder we compare.

\GOneDetail{Year coverage}
\GOneLabel{sec:wrdg-year-coverage}
For production in year $t$, we obtain beginning capital from the stock reported in calendar year $t-1$, treating that amount as its year-end stock. Materials or investment comes from year $t$ itself. The dynamic stage also needs the preceding production observation; its beginning capital consequently draws on year $t-2$. We preserve every calendar gap instead of treating the next observed year as the next calendar year. Table~\GOneRef{tab:deep-wrdg-input-states} reports each input's availability before we require them together.
\GOneInput{tables/deep_wrdg_input_states.tex}
The positive observations for one input need not be the same observations as those for another. We retain a record only when it supplies all the quantities that its estimation step requires. A missing response does not show zero production, and we cannot take the log of a recorded zero. Beginning capital also requires the immediately preceding stock. The year-end and investment assumptions admit positive reported amounts without changing the documentary flags; those flags continue to show where the source evidence is narrower than the assumption.

\GOneDetail{Transformations and sample restrictions}
\GOneLabel{sec:wrdg-transformations}
We check positivity before taking logs. We use the available production years and sectors, match each record to the immediately preceding calendar year, and apply the minimum sample sizes to the pooled sector sample. For the polynomial calculations, we then subtract the relevant sample median from each logged variable; this changes its numerical reference point while keeping the eligible records fixed. We estimate one coefficient vector and control relationship per sector over the pooled years, without adding year effects. Table~\GOneRef{tab:deep-wrdg-sample} traces the sample by sector and year.
\GOneInput{tables/deep_wrdg_sample.tex}
The estimation steps use these records in different ways. An earlier observation can supply the preceding-year inputs for a later calendar pair. The estimators with a separate first stage can also use eligible records outside the accepted pairs in that first-stage calculation. Once a sector has estimated coefficients, we calculate a residual for each record with the required current inputs, including one without an eligible predecessor. Calculating that residual does not add a new pair or re-estimate the coefficients. An observation outside the dynamic-stage sample can therefore contribute to the first stage, receive a residual, or do both.


\GOneSubsubsection{System equations}
\GOneLabel{sec:wrdg-economics}
Writing these two explanations of logged value added gives
\begin{equation}
\begin{aligned}
y_{it}&=\alpha+\beta_Ll_{it}+\beta_Kk_{it}+h(k_{it},m_{it})+\varepsilon_{it},\\
y_{it}&=\alpha+\beta_Ll_{it}+\beta_Kk_{it}+g\!\left[h(k_{i,t-1},m_{i,t-1})\right]+\xi_{it}+\varepsilon_{it}.
\end{aligned}
\GOneLabel{eq:wrdg-general}
\end{equation}
The intercept $\alpha$ captures a constant component common to these expressions. Current labor can respond to $\xi_{it}$, so it may be unrelated to the first equation's error while still related to the second equation's error. We therefore use current labor among the first equation's instruments, and previous-year labor with beginning-of-year capital for the second. Estimating the equations jointly uses both sets of restrictions; the first equation need not identify each input coefficient separately on its own.

\GOneSubsubsection{Productivity transition}
\GOneLabel{sec:wrdg-restriction}
We impose a specific relationship between current and previous-year productivity. We use the same coefficients on the current and previous-year proxy polynomials. Within this approximation, we predict current productivity as previous productivity plus a constant change. Thus an advantage in the previous year carries forward without shrinking or growing; a new surprise changes its subsequent level. This is a random-walk specification with a drift, the constant expected change. Wooldridge describes this special case in his author lecture \citep[pp.~21--22]{wooldridgelecture}. It restricts the general function $g(\cdot)$ in Equation~\GOneEqref{eq:wrdg-general}.

To implement that restriction, we collect the labor and capital coefficients, the common intercept $a_0$, the drift $e_0$ and the polynomial coefficients $\theta$ in $\vartheta=(b_L,b_K,a_0,e_0,\theta')'$. Subtracting each equation's predicted value from observed logged value added gives two errors:
\begin{equation}
\begin{aligned}
u_{1,it}(\vartheta)&=y_{it}-b_Ll_{it}-b_Kk_{it}-a_0-B_3(k_{it},m_{it})'\theta,\\
u_{2,it}(\vartheta)&=y_{it}-b_Ll_{it}-b_Kk_{it}-a_0-e_0-B_3(k_{i,t-1},m_{i,t-1})'\theta.
\end{aligned}
\GOneLabel{eq:wooldridge}
\end{equation}
The second equation includes the drift $e_0$, while both equations use the same $\theta$. We therefore estimate a shared relationship for current and previous-year proxy values. LP and ACF instead refit a separate cubic relationship between previous-year and current productivity. This difference in the assumed evolution of productivity is part of the Wooldridge specification reported here.

In the quadratic sensitivity specification, we replace both cubic proxy polynomials, $B_3$, with quadratic polynomials, $B_2$, and retain the shared coefficients. This changes the proxy approximation while preserving the random-walk restriction on productivity.

\GOneSubsubsection{Estimation procedure}
\GOneLabel{sec:wrdg-stages}
We choose instruments separately for the two errors. For the first equation, we use an intercept, current labor and the current capital--materials polynomial. For the second, we use an intercept, beginning-of-year capital, previous-year labor and the previous-year capital--materials polynomial. These choices give
\begin{equation}
\begin{aligned}
z_{1,it}&=(1,l_{it},B_3(k_{it},m_{it})')',\\
z_{2,it}&=(1,k_{it},l_{i,t-1},B_3(k_{i,t-1},m_{i,t-1})')'.
\end{aligned}
\GOneLabel{eq:wrdg-instruments}
\end{equation}
For each consecutive-year observation, we multiply each equation's error by each of its instruments. We combine these products, $z_1u_1$ and $z_2u_2$, in $v_{it}(\vartheta)$ and average them across the common estimation sample, $\mathcal D$. As in ACF, the assumed population averages are zero. We choose the coefficients to minimize a weighted measure of the sample averages' departures from zero:
\begin{equation}
Q_W(\vartheta)=\overline v(\vartheta)'W\overline v(\vartheta),\qquad
\overline v(\vartheta)=N^{-1}\sum_{(i,t)\in\mathcal D}v_{it}(\vartheta).
\GOneLabel{eq:wrdg-gmm}
\end{equation}
The matrix $W$ determines the relative weight given to these departures. We estimate the system twice: first using an independent, unadjusted weight matrix, then updating the weights using the variation in the fitted moment products. We solve the minimization with Gauss--Newton optimization, which uses local changes in the criterion to update trial coefficients. We also check the independent information in the instruments and whether changes in the coefficients produce distinguishable changes in the moments. These instrument and moment-derivative ranks address different questions from the rank of the estimated coefficient covariance matrix. Both equations use the same consecutive-year sample.

\GOneSubsubsection{Worked aggregate calculation}
\GOneLabel{sec:wrdg-worked}
We illustrate the calculation using mean logged inputs within one sector. Multiplying each mean log input by its fitted coefficient and subtracting both contributions from mean logged value added gives the mean individual residual exactly, because the same slopes apply to every record in that sector. Table~\GOneRef{tab:deep-wrdg-worked} follows Equations~\GOneEqref{eq:native-definition}--\GOneEqref{eq:relative-index} through centering and exponentiation, using full precision before rounding.
\GOneInput{tables/deep_wrdg_worked.tex}
In 2011, mean logged value added is 13.7219. The fitted labor and capital contributions are 2.2691 and 1.2040, leaving 10.2488. Subtracting the pooled reference mean of 10.9397 gives a centered mean of -0.6910. Its exponential is 0.5011, while the arithmetic mean of individual index values is 2.0340. The latter gives greater weight to large individual residuals after exponentiation. Both annual examples use the same pooled equation; neither estimates a separate annual technology.


\GOneSubsubsection{Coefficient estimates}
\GOneLabel{sec:wrdg-coefficients}
We record the primary outcome for every sector and both input-sample definitions in Table~\GOneRef{tab:deep-wrdg-attempts}. This keeps the absence of a fit separate from a coefficient estimated near zero.
\GOneInput{tables/deep_wrdg_attempts.tex}
Only estimated cells supply coefficients and derived residuals. A count requirement and a numerical failure have different meanings: the former prevents the declared calculation from being attempted, while the latter prevents an attempted calculation from supplying an accepted estimate.
We report the estimated coefficients beginning with Table~\GOneRef{tab:deep-wrdg-regression}. The own and common samples have separate columns and fits. To describe uncertainty, we resample whole fictional company histories and repeat every estimation step. Standard errors measure the spread of these repeated estimates; percentile intervals give their admitted lower and upper quantiles.
\GOneInput{tables/deep_wrdg_regression.tex}
\GOneInput{tables/deep_wrdg_regression-2.tex}
\GOneInput{tables/deep_wrdg_regression-3.tex}
\GOneInput{tables/deep_wrdg_regression-common.tex}
The uncertainty in the coefficient sum depends on how labor and capital estimates move together across resamples. Adding their separate standard errors would ignore that covariance. Fitted year coverage is pooled, and a separate first stage can use more records than the dynamic sample. These counts describe distinct roles in the same estimation.

Table~\GOneRef{tab:deep-wrdg-numerics} reports the rank and optimization diagnostics for the successful primary fits. We compare numerical rank with the dimension of the relevant matrix because returning a coefficient is not enough to show that the stated calculation had the required independent information.
\GOneInput{tables/deep_wrdg_numerics.tex}
The diagnostics describe the numerical problem at the reported solution. They do not establish valid economic timing or a unique global optimum; the alternative-start comparison examines a separate aspect of that uncertainty.

Figure~\GOneRef{fig:deep-wrdg-coefficients} displays the same own-sample coefficients and full admitted uncertainty intervals. Separate input panels preserve the scale of the labor and capital estimates.
\begin{figure}[H]
\centering
\caption{Wooldridge: coefficient estimates}
\GOneLabel{fig:deep-wrdg-coefficients}
\includegraphics[width=\linewidth]{figures/deep_wrdg_coefficients.pdf}
\appnotes{We pool supported years within each sector under constant coefficients. Marks show own-sample primary estimates and intervals admitted under the bootstrap minimum-success rule. Intervals retain their full endpoints and condition on the fictional histories and maintained measurements.}{G0.5 v3 modified manufacturing records with fictional company links; authors' calculations.}
\end{figure}
Across the estimated sectors, labor ranges from 0.1437 to 0.8631; capital ranges from -0.0169 to 0.1038. We retain the signs and magnitudes returned by the primary specification. Their interpretation depends on beginning-capital timing, reported monetary scales, the proxy assumptions and the estimation criterion. The intervals describe additional variation when the same calculation is repeated on resampled histories. The capital slope is negative in sectors 14 and 32. At given reported value added and workers, higher beginning capital therefore raises the calculated residual in those equations. This sign is a feature of the nominal fitted relationship; it does not establish that additional productive capital reduces actual output. It changes how capital differences enter the residual ranking and must be considered when interpreting that ordering.
\GOneSubsubsection{Bootstrap coefficient distributions}
\GOneLabel{sec:wrdg-bootstrap}
Table~\GOneRef{tab:deep-wrdg-bootstrap} reports the spread of the repeated coefficient estimates. The standard error is their sample standard deviation, while the median and percentiles reveal features that a single standard error can conceal. The minimum and maximum make finite extremes visible.
\GOneInput{tables/deep_wrdg_bootstrap.tex}
The largest labor-coefficient standard deviation occurs in sector 25, at 0.1464. Its bootstrap median is 0.6931, with P5 and P95 at 0.4166 and 0.8802; the full range is 0.2803 to 1.0434. Comparing the middle of this distribution with its extremes helps distinguish broad dispersion from a small number of distant estimates. We retain every finite resampled estimate and keep the original-sample coefficient as the primary result. A limited number of resamples determines the extreme percentiles, so their locations should not be interpreted as precise tail probabilities. The exercise also leaves uncertainty about actual identities, prices and economic restrictions outside the resampling calculation.
\GOneSubsubsection{Residual distributions}
\GOneLabel{sec:wrdg-distribution}
We subtract the fitted labor and capital contributions from logged value added, then subtract each sector's pooled dynamic residual mean. Figure~\GOneRef{fig:deep-wrdg-histogram} compares the uncentered and centered distributions on exactly the same fitted records.
\begin{figure}[H]
\centering
\caption{Wooldridge: pooled residual distributions}
\GOneLabel{fig:deep-wrdg-histogram}
\includegraphics[width=\linewidth]{figures/deep_wrdg_histogram.pdf}
\appnotes{Both panels contain the same own-sample dynamic observations pooled across all fitted years. Uncentered residuals use original logs; centering subtracts one mean per sector-method-sample. Bins span 0.5 log points and retain all finite tails. Labels show the modal and first and last occupied bins with source-derived counts and shares.}{G0.5 v3 modified manufacturing records with fictional company links; authors' calculations.}
\end{figure}
The centered distribution has a standard deviation of 1.422 log points. Centering removes differences in sector means while preserving all within-sector gaps and rankings. The pooled shape can change because it previously combined those mean differences with variation within sectors. We retain all finite observations, including the tails.
\GOneSubsubsection{Prediction samples and tails}
\GOneLabel{sec:wrdg-prediction}
The histogram above describes the dynamic estimation sample. We can calculate residuals more broadly whenever a record supplies the required current inputs in a sector with estimated slopes. Table~\GOneRef{tab:deep-wrdg-scopes} distinguishes these roles for every source year and for the pooled records.
\GOneInput{tables/deep_wrdg_scopes.tex}
There are 10,315 predicted company-year observations: 5,041 belong to the pooled dynamic fit and 5,274 lie outside it. Only the dynamic-fit records determine the reference mean. The annual fit means and the broader prediction means can therefore differ from zero, even though the pooled fit mean is zero by construction. A company appearing repeatedly contributes several observations to these counts.

Table~\GOneRef{tab:deep-wrdg-quantiles} holds the fitted records fixed and compares the uncentered logs, centered logs and positive index. We include the minimum and maximum so that the transformed tails remain visible.
\GOneInput{tables/deep_wrdg_quantiles.tex}
Centering preserves every within-sector gap. Exponentiation preserves within-sector ordering but enlarges distances at high values, as Equation~\GOneEqref{eq:relative-index} shows. A large index maximum therefore needs to be read through this transformation, rather than treated directly as physical efficiency.

Table~\GOneRef{tab:deep-wrdg-tails} identifies distant observations using two distinct rules. The interquartile range (IQR) is the distance from P25 to P75. Fences placed 1.5 IQR beyond those quartiles compare a value with the spread of its own distribution. Fixed absolute-log and index cutoffs instead apply the same numerical threshold in every sector. A record can exceed more than one cutoff.
\GOneInput{tables/deep_wrdg_tails.tex}
In the pooled dynamic sample, 679 records have an absolute centered log residual above two. These flags describe the residuals produced by the fitted equation. Prices, reported amounts, the estimated input contributions and production conditions can all affect an extreme residual, so the flag alone identifies neither an error nor its cause. We retain every flagged observation.

\GOneSubsubsection{Annual distributions}
\GOneLabel{sec:wrdg-annual}
Table~\GOneRef{tab:deep-wrdg-annual} compares all eligible predictions by sector and year. We keep each sector's fitted slopes and reference mean constant over time. This preserves annual differences instead of subtracting each year's average separately.
\GOneInput{tables/deep_wrdg_annual.tex}
Changes in annual means and quartiles combine differences in reported outcomes and inputs with differences in sample composition. A year with more high-residual observations can have a higher mean without every continuing assignment improving. We examine this distinction next by matching observations within adjacent years.
Figure~\GOneRef{fig:deep-wrdg-year-sector} and its continuation panels plot the annual median and quartiles for each sector. These intervals measure dispersion across observations, not uncertainty in an estimated coefficient.
\begin{figure}[H]
\centering
\caption{Wooldridge annual residual quartiles}
\GOneLabel{fig:deep-wrdg-year-sector}
\includegraphics[width=\linewidth]{figures/deep_wrdg_year_sector.pdf}
\appnotes{Each sector panel includes the available annual predictions, using one pooled coefficient vector and reference mean. Dots mark medians and intervals run from P25 to P75; all available values are numerically labelled. An empty group is not a zero. The legend identifies any missing-file region. Continuation panels use the same construction.}{G0.5 v3 modified manufacturing records with fictional company links; authors' calculations.}
\end{figure}
\begin{figure}[H]
\centering
\caption{Wooldridge annual residual quartiles (continued)}
\GOneLabel{fig:deep-wrdg-year-sector-p2}
\includegraphics[width=\linewidth]{figures/deep_wrdg_year_sector_p2.pdf}
\appnotes{Each sector panel includes the available annual predictions, using one pooled coefficient vector and reference mean. Dots mark medians and intervals run from P25 to P75; all available values are numerically labelled. An empty group is not a zero. The legend identifies any missing-file region. Continuation panels use the same construction.}{G0.5 v3 modified manufacturing records with fictional company links; authors' calculations.}
\end{figure}
\begin{figure}[H]
\centering
\caption{Wooldridge annual residual quartiles (continued)}
\GOneLabel{fig:deep-wrdg-year-sector-p3}
\includegraphics[width=\linewidth]{figures/deep_wrdg_year_sector_p3.pdf}
\appnotes{Each sector panel includes the available annual predictions, using one pooled coefficient vector and reference mean. Dots mark medians and intervals run from P25 to P75; all available values are numerically labelled. An empty group is not a zero. The legend identifies any missing-file region. Continuation panels use the same construction.}{G0.5 v3 modified manufacturing records with fictional company links; authors' calculations.}
\end{figure}

The sector panels separate variation within each industry from changes in which industries dominate a pooled annual distribution. They describe the same fixed sector equations over time.
\GOneSubsubsection{Adjacent-year changes}
\GOneLabel{sec:wrdg-paired}
We hold membership fixed within each comparison by requiring a prediction for the same fictional assignment and sector in adjacent calendar years. Table~\GOneRef{tab:deep-wrdg-pair-coverage} shows how this restriction narrows the endpoint samples. We retain intervals with no matched predictions so that missing capital history remains visible.
\GOneInput{tables/deep_wrdg_pair_coverage.tex}
The pooled comparison contains 5,041 record-pair intervals. A company can enter several successive intervals, so this is not a count of independent companies. Restricting to matched predictions removes changes in membership within each interval, although membership can still differ between intervals.

For each pair, we subtract the earlier centered residual from the later one. Let $\Delta x_{it}=x_{it}-x_{i,t-1}$ for any logged variable $x$. Holding the sector slopes fixed gives
\begin{equation}
\Delta\widetilde r_{it}=\Delta y_{it}-\widehat\beta_{L,s}\Delta l_{it}-\widehat\beta_{K,s}\Delta k_{it}.
\GOneLabel{eq:wrdg-paired-change}
\end{equation}
The pooled reference mean cancels because we subtract the same constant in both years. Table~\GOneRef{tab:deep-wrdg-paired} reports the resulting changes for each interval and sector, then pools the record-pair intervals under the same rule.
\GOneInput{tables/deep_wrdg_paired.tex}
The pooled mean log change is 0.146. The correlation of earlier and later ranks is 0.775, measuring persistence in ordering. The mean individual index ratio is 1.956; large ratios can strongly influence this arithmetic mean, which need not equal the exponential of the mean log change. These comparisons use fictional links and reported-price equations; they do not establish real-company productivity changes.
Figure~\GOneRef{fig:deep-wrdg-paired-change} and its continuation panels plot these adjacent-year changes under the same matched-prediction rule.
\begin{figure}[H]
\centering
\caption{Wooldridge adjacent-year residual changes}
\GOneLabel{fig:deep-wrdg-paired-change}
\includegraphics[width=\linewidth]{figures/deep_wrdg_paired_change.pdf}
\appnotes{Each comparison uses same-sector fictional assignments with eligible predictions in consecutive calendar years. Sector coefficients and pooled reference means remain fixed. Histograms retain all finite changes. Numeric labels identify the modal and first and last occupied bins. Pooled panels combine record-pair intervals, allowing repeated company contributions; interval panels retain their own denominators.}{G0.5 v3 modified manufacturing records with fictional company links; authors' calculations.}
\end{figure}
\begin{figure}[H]
\centering
\caption{Wooldridge adjacent-year residual changes (continued)}
\GOneLabel{fig:deep-wrdg-paired-change-p2}
\includegraphics[width=\linewidth]{figures/deep_wrdg_paired_change_p2.pdf}
\appnotes{Each comparison uses same-sector fictional assignments with eligible predictions in consecutive calendar years. Sector coefficients and pooled reference means remain fixed. Histograms retain all finite changes. Numeric labels identify the modal and first and last occupied bins. Pooled panels combine record-pair intervals, allowing repeated company contributions; interval panels retain their own denominators.}{G0.5 v3 modified manufacturing records with fictional company links; authors' calculations.}
\end{figure}
\begin{figure}[H]
\centering
\caption{Wooldridge adjacent-year residual changes (continued)}
\GOneLabel{fig:deep-wrdg-paired-change-p3}
\includegraphics[width=\linewidth]{figures/deep_wrdg_paired_change_p3.pdf}
\appnotes{Each comparison uses same-sector fictional assignments with eligible predictions in consecutive calendar years. Sector coefficients and pooled reference means remain fixed. Histograms retain all finite changes. Numeric labels identify the modal and first and last occupied bins. Pooled panels combine record-pair intervals, allowing repeated company contributions; interval panels retain their own denominators.}{G0.5 v3 modified manufacturing records with fictional company links; authors' calculations.}
\end{figure}
\begin{figure}[H]
\centering
\caption{Wooldridge adjacent-year residual changes (continued)}
\GOneLabel{fig:deep-wrdg-paired-change-p4}
\includegraphics[width=\linewidth]{figures/deep_wrdg_paired_change_p4.pdf}
\appnotes{Each comparison uses same-sector fictional assignments with eligible predictions in consecutive calendar years. Sector coefficients and pooled reference means remain fixed. Histograms retain all finite changes. Numeric labels identify the modal and first and last occupied bins. Pooled panels combine record-pair intervals, allowing repeated company contributions; interval panels retain their own denominators.}{G0.5 v3 modified manufacturing records with fictional company links; authors' calculations.}
\end{figure}

A pooled average can conceal increases and decreases across matched assignments. The interval panels show how their distributions vary over time under the fixed sector equations.
\GOneSubsubsection{Specification sensitivity}
\GOneLabel{sec:wrdg-sensitivity}
Figure~\GOneRef{fig:deep-wrdg-sensitivity} and its continuation panels compare the prespecified starting values and proxy-polynomial alternatives. Changing the start leaves the model fixed; changing the polynomial alters the fitted control relationship. A numerical stopping rule need not identify the same solution from every starting value.
\begin{figure}[H]
\centering
\caption{Wooldridge: estimation sensitivity}
\GOneLabel{fig:deep-wrdg-sensitivity}
\includegraphics[width=\linewidth]{figures/deep_wrdg_sensitivity.pdf}
\appnotes{The primary cubic specification, two alternative starts and quadratic proxy use the same own-sample eligibility and numerical translations. Objective values are comparable only within a method and matching polynomial specification. Missing alternatives remain visible as missing; continuation panels retain identical construction.}{G0.5 v3 modified manufacturing records with fictional company links; authors' calculations.}
\end{figure}
\begin{figure}[H]
\centering
\caption{Wooldridge: estimation sensitivity (continued)}
\GOneLabel{fig:deep-wrdg-sensitivity-p2}
\includegraphics[width=\linewidth]{figures/deep_wrdg_sensitivity_p2.pdf}
\appnotes{The primary cubic specification, two alternative starts and quadratic proxy use the same own-sample eligibility and numerical translations. Objective values are comparable only within a method and matching polynomial specification. Missing alternatives remain visible as missing; continuation panels retain identical construction.}{G0.5 v3 modified manufacturing records with fictional company links; authors' calculations.}
\end{figure}
Alternative starts change the labor coefficient by at most 1.09e-14; alternative starts change the capital coefficient by at most 1.70e-15; the quadratic proxy change the labor coefficient by at most 4.64e-02; the quadratic proxy change the capital coefficient by at most 3.72e-02. These comparisons keep each sector's records fixed. We do not use an economically convenient sign, a narrow residual distribution or a smaller criterion under a different model to select a replacement primary estimate.



\GOneSubsection{Coefficient estimates and uncertainty}
\GOneLabel{sec:coefficients}
The method sections report separate estimates for each sector. We now bring together 37 primary own-sample estimates and 4 primary common-sample estimates to compare the contributions assigned to labor and capital. For each model, we also repeat the estimation on samples formed by drawing complete fictional company histories with replacement within the sector. Drawing whole histories keeps each sampled company's consecutive observations together. This bootstrap calculation shows how the estimated coefficients vary when the composition of the sample changes under the same fictional links and measurement assumptions.

Each model requests 99 resamples. The own-sample estimates return 3,663 successful resamples, and the common-sample estimates return 396. The tables retain the successful count for each model. We summarize the resulting estimates with a standard error, calculated from their spread, and an exploratory percentile interval that spans the middle 95\% of their distribution. We report an interval only when at least 95 resamples succeed. Table~\GOneRef{tab:coefficients} places the coefficient estimates beside these intervals; the method-specific tables also report the standard errors. Resampling the available histories does not capture uncertainty about whether the fictional links match actual companies, nor does it correct the reported monetary values to common prices.
\GOneInput{tables/coefficients.tex}
Across own-sample estimates, the labor coefficient ranges from 0.1437 (Wooldridge, sector 12) to 1.4197 (ACF, sector 12); the capital coefficient ranges from -0.0209 (LP, sector 22) to 0.1735 (ACF, sector 12). These ranges describe variation across both methods and sectors. They do not isolate a method effect unless we also fix the sector and comparison sample. We retain the estimated signs rather than constraining the coefficients after seeing them. The largest labor standard error is 3.8544, for ACF in sector 15; its admitted percentile interval is [1.0290, 16.5219]. The standard error uses all successful resampled estimates, while the percentile interval excludes their outer tails.

The coefficient sum measures the model's response when both inputs rise proportionally. A physical returns-to-scale interpretation requires comparable production quantities, which these nominal measures do not establish.

Figure~\GOneRef{fig:coefficients} shows the own-sample estimates and full admitted intervals from resampled histories. The table also gives the separate common-sample estimates.
\begin{figure}[H]
\centering
\caption{Own-sample input coefficients and uncertainty}
\GOneLabel{fig:coefficients}
\includegraphics[width=\linewidth]{figures/coefficients.pdf}
\appnotes{Marks show primary own-sample coefficients. Available percentile intervals use successful full-history bootstrap resamples and the declared minimum-success rule. We retain full interval endpoints; the table separately includes common-sample estimates.}{G0.5 v3 modified manufacturing records with fictional company links; authors' calculations.}
\end{figure}
We read the coefficient positions together with the interval widths. A large difference between methods and an imprecise estimate answer different questions: the first concerns the equation assigned to the records, while the second concerns variation across resampled histories. The residual distributions show how these estimated equations translate into differences between observations.

\GOneSubsection{Residual distributions}
\GOneLabel{sec:distributions}
We apply each method's sector coefficients to the corresponding logged value added, workers and beginning capital. We subtract the method-sector mean over all of its pooled dynamic estimation records, following Section~\GOneRef{sec:meaning}:
\begin{equation}
\widehat r_{it}^{(m)}=y_{it}-\widehat\beta_{L,s}^{(m)}l_{it}-\widehat\beta_{K,s}^{(m)}k_{it},\qquad
\widetilde r_{it}^{(m)}=\widehat r_{it}^{(m)}-\frac{1}{N_s^{(m)}}\sum_{(j,\tau)\in\mathcal D_s^{(m)}}\widehat r_{j\tau}^{(m)}.
\GOneLabel{eq:residual}
\end{equation}
The summation includes all fitted company-year records in that sector and method, using $\tau$ to distinguish their years from the prediction year $t$. Table~\GOneRef{tab:distribution} reports the spread and selected percentiles for the own and common estimation samples. A percentile marks the value below which the corresponding share of observations falls; the median is the 50th percentile. Separate own samples can differ in their records, so differences in their pooled distributions can reflect both estimation and composition.
\GOneInput{tables/distribution.tex}
The standard deviations of own-sample centered residuals are OP 1.363, LP 1.467, ACF 1.130, Wooldridge 1.422 log points. A narrower distribution leaves less variation after the fitted input contributions, but does not establish better measurement of physical productivity. These own-sample summaries also differ in sector and observation coverage. The common-sample rows and matched comparisons help separate that coverage issue from differences between fitted equations.

Figure~\GOneRef{fig:distribution} groups the residuals into equal-width intervals and plots the share of observations in each interval. These histograms show where the records concentrate and how far the less frequent values extend. We retain all occupied intervals, including the tails.
\begin{figure}[H]
\centering
\caption{Own-sample centered residual distributions}
\GOneLabel{fig:distribution}
\includegraphics[width=\linewidth]{figures/distribution.pdf}
\appnotes{We pool dynamic estimation records after subtracting each sector-method-sample fitted mean. Bins span 0.5 log points and retain all finite tails. Numeric labels identify the modal and first and last occupied bins. Each method uses its own denominator.}{G0.5 v3 modified manufacturing records with fictional company links; authors' calculations.}
\end{figure}
Subtracting each sector's average gives the methods a common zero mean, but leaves their medians, spread and tails free to differ. The histograms and percentiles therefore describe differences that centering does not remove. Similar distribution shapes still need not place individual records in the same order, which we examine next.

Figure~\GOneRef{fig:deep-common-distribution} compares the four methods after requiring inputs suitable for all of them. The common intersection contains 128 fitted company-year records. Each model is estimated separately under the common-input rule, so this fixes the observations in the comparison without forcing the methods to use identical coefficients.
\begin{figure}[H]
\centering
\caption{Common-sample centered residual distributions}
\GOneLabel{fig:deep-common-distribution}
\includegraphics[width=\linewidth]{figures/deep_common_distribution.pdf}
\appnotes{The four panels use common-input primary models on their shared pooled dynamic company-year records. Each method has separately estimated coefficients and its own pooled sector-method reference mean. Bins span 0.5 log points and use the same display axes across methods. Numeric labels identify the modal and first and last occupied bins. Every finite tail remains; the common sample is distinct from the wider own-input samples.}{G0.5 v3 modified manufacturing records with fictional company links; authors' calculations.}
\end{figure}
The common-sample standard deviations are OP 1.362, LP 1.464, ACF 1.222, Wooldridge 1.453 log points. Since these panels use the same records, their differences describe how the fitted methods allocate value added between input contributions and the remainder. These common-sample distributions do not replace the wider own-sample results: requiring inputs suitable for all four methods changes the records available for estimation and can change an earlier regression even when a dynamic pair remains the same.


\GOneSubsection{Comparison across estimators}
\GOneLabel{sec:agreement}
We match each method's residual to the same source observation before comparing methods. We use two measures of agreement. Pearson correlation describes how closely the residual values vary together in a straight-line relationship. Spearman correlation applies that comparison to ranks, asking whether the methods place the same observations high or low. A Pearson correlation of one means that the residual values lie on an increasing straight line; it does not mean that the values are equal. A Spearman correlation of one means that the rankings agree. Zero indicates no linear association between the values, or between their ranks in the Spearman calculation; a negative correlation indicates an opposite association.

We calculate both correlations after subtracting the sector-specific means. Table~\GOneRef{tab:rank-comparison} reports the six possible method pairs under both the own and common input-sample specifications. Its sample columns distinguish records with predictions from both methods in a pair from records with predictions from all four methods.
\GOneInput{tables/rank_comparison.tex}
The own-sample rank correlations are OP--LP: 0.966 on 128 shared records; OP--ACF: 0.949 on 128 shared records; OP--Wooldridge: 0.974 on 128 shared records; LP--ACF: 0.845 on 5,041 shared records; LP--Wooldridge: 0.997 on 5,041 shared records; ACF--Wooldridge: 0.866 on 5,041 shared records. Each pair holds its records fixed, but different pairs can use different intersections. The common-sample rows compare the separately fitted common-input specifications; we do not equate those coefficients with the own-sample coefficients. Correlation measures agreement, not which estimator best measures physical productivity.

Figure~\GOneRef{fig:rank-comparison} repeats the rank comparisons separately within each sector. This shows whether agreement in the pooled sample also holds among records in the same industry group.
\begin{figure}[H]
\centering
\caption{Own-sample rank correlations by sector}
\GOneLabel{fig:rank-comparison}
\includegraphics[width=\linewidth]{figures/rank_comparison.pdf}
\appnotes{Each coefficient compares the same original company-year records in a method pair. Pooled correlations combine available sectors after centering. An unavailable pair is not zero agreement. Different pairs can use different intersections.}{G0.5 v3 modified manufacturing records with fictional company links; authors' calculations.}
\end{figure}
Each comparison now fixes both the records and their sector. Differences across the plotted correlations indicate that agreement between methods can vary across industry groups. The pooled coefficient consequently cannot describe every sector's agreement by itself.

We also examine whether agreement differs between low- and high-ranked observations. For each method pair, we order records by the first method's residual rank and divide those ranks into groups. Figure~\GOneRef{fig:rank-scatter} plots the paired averages within each group. We omit a group if it contains fewer than five records.
\begin{figure}[H]
\centering
\caption{Binned residual comparisons}
\GOneLabel{fig:rank-scatter}
\includegraphics[width=\linewidth]{figures/rank_scatter.pdf}
\appnotes{We rank the first method's centered residual on shared dynamic records, assign average ranks to ties and form 20 bins. We show a bin only if it contains at least five records. Each mark summarizes paired residuals or ranks as indicated on its axes; sparse bins remain outside the plotted summaries.}{G0.5 v3 modified manufacturing records with fictional company links; authors' calculations.}
\end{figure}
Each plotted point summarizes several records. The grouping helps show broad differences along the ranking, but hides differences within each group; omitted small groups also remain outside the figure. The full-sample correlations in Table~\GOneRef{tab:rank-comparison} include those observations. We therefore use the grouped plots to describe where agreement changes, alongside the correlations calculated from all shared records.

\GOneSubsection{Specification sensitivity}
\GOneLabel{sec:sensitivity}
The estimates can change even when we keep the observations fixed. A numerical search begins from specified coefficient values, and a nonlinear search may end at different solutions when it starts elsewhere. The flexibility of the functions used to represent unobserved production conditions is a separate estimation choice. We examine both by repeating each sector--method calculation with two fixed alternative starting values and with a quadratic, rather than cubic, proxy approximation. The latter removes the third-degree terms from the proxy function described in Section~\GOneRef{sec:common-basis}.

Across the own and common input samples, the primary and alternative specifications produce 164 sets of point estimates. Within each paired comparison we hold the sample fixed. Table~\GOneRef{tab:sensitivity} summarizes the largest absolute changes across sectors for the own-sample estimates. Separate columns distinguish changing where the search starts from changing the proxy approximation.
\GOneInput{tables/sensitivity.tex}
The largest absolute change under alternative starts is 0.9262, for ACF's labor coefficient. The quadratic proxy comparison reaches 0.1818, for OP's capital coefficient. These maxima identify sensitivity within the tested choices, keeping each compared method-sector sample fixed. Changing the approximation changes the fitted model as well as the numerical problem.

Figure~\GOneRef{fig:sensitivity} identifies which sector--method estimates account for the changes in the labor coefficient. We display the alternatives alongside the primary estimate, including alternatives with less conventional signs or magnitudes.
\begin{figure}[H]
\centering
\caption{Labor-coefficient sensitivity}
\GOneLabel{fig:sensitivity}
\includegraphics[width=\linewidth]{figures/sensitivity.pdf}
\appnotes{We keep the own-sample records fixed within each sector and method. Alternative starts change only initial structural coefficients. The quadratic specification changes the proxy approximation. The primary estimate remains the one specified before examining these comparisons.}{G0.5 v3 modified manufacturing records with fictional company links; authors' calculations.}
\end{figure}
We retain the primary specifications selected before examining these alternatives. Reaching a numerical stopping rule establishes neither a unique optimum nor the validity of the economic assumptions. The alternative coefficients show how much the assigned input contributions depend on the particular search and approximation choices we examine.

\GOneSubsection{Conclusion}
\GOneLabel{sec:synthesis}
Treating the reported capital totals as year-end stocks permits 37 own-sample sector estimates and 4 common-sample estimates. The materials-based methods cover more of the available records, while the comparison shared by all four methods contains 128 fitted company-year observations. We therefore distinguish broad prediction coverage from the narrower basis for a four-method comparison.

LP and Wooldridge give closely aligned rankings on 5,041 shared own-sample records, with a rank correlation of 0.997. The LP--ACF correlation is 0.845. These comparisons show that methods can agree strongly about ordering without assigning the same contributions to the inputs. Uncertainty also differs across equations: the largest labor bootstrap standard error is 3.8544, for ACF in sector 15. Keeping the records and model fixed but changing the numerical start moves ACF's labor coefficient by as much as 0.9262. Agreement about residual ordering alone cannot establish stable input contributions or validate the assumed production relationship.

These findings describe calculations using modified monetary records and fictional company histories. The all-years stock assumption makes it possible to account for beginning capital wherever the immediately preceding stock exists. It does not authenticate the meaning or scale of every year's capital field. Pooling the years further assumes stable sector coefficients and control relationships; we preserve the missing calendar year and use a single pooled reference when comparing annual residuals.

Interpreting the remainder as productivity of actual manufacturers requires verified company links, comparable production and input measures at common prices, and defensible assumptions about input choices and continuation. The calculations here show what follows under the stated assumptions and where uncertainty or sample requirements limit the comparison.

\begin{thebibliography}{9}
\bibitem[Ackerberg, Caves, and Frazer(2015)]{acf} Ackerberg, Daniel A., Kevin Caves, and Garth Frazer. 2015. ``Identification Properties of Recent Production Function Estimators.'' \textit{Econometrica} 83(6): 2411--2451. \href{https://doi.org/10.3982/ECTA13408}{doi:10.3982/ECTA13408}.
\bibitem[Levinsohn and Petrin(2003)]{lp} Levinsohn, James, and Amil Petrin. 2003. ``Estimating Production Functions Using Inputs to Control for Unobservables.'' \textit{Review of Economic Studies} 70(2): 317--341. \href{https://www.nber.org/papers/w7819}{Working-paper and published-version record}.
\bibitem[Olley and Pakes(1996)]{op} Olley, G. Steven, and Ariel Pakes. 1996. ``The Dynamics of Productivity in the Telecommunications Equipment Industry.'' \textit{Econometrica} 64(6): 1263--1297. \href{https://www.nber.org/papers/w3977}{Working-paper and published-version record}.
\bibitem[Rovigatti and Mollisi(2018)]{prodest} Rovigatti, Gabriele, and Vincenzo Mollisi. 2018. ``Theory and Practice of Total-Factor Productivity Estimation: The Control Function Approach Using Stata.'' \textit{Stata Journal} 18(3): 618--662. \href{https://www.stata-journal.com/article.html?article=st0537}{Article and implementation record}.
\bibitem[Wooldridge(2009)]{wooldridge} Wooldridge, Jeffrey M. 2009. ``On Estimating Firm-Level Production Functions Using Proxy Variables to Control for Unobservables.'' \textit{Economics Letters} 104(3): 112--114. \href{https://doi.org/10.1016/j.econlet.2009.04.026}{doi:10.1016/j.econlet.2009.04.026}.
\bibitem[Wooldridge(2009b)]{wooldridgelecture} Wooldridge, Jeffrey M. 2009b. ``New Developments in Econometrics: Lecture 4, Linear Panel Data Models II.'' Cemmap/UCL course materials, June. \href{https://www.cemmap.ac.uk/wp-content/legacy/resources/imbens_wooldridge/slides_4.pdf}{Author lecture; shared-vector special case, printed pp.~21--22}.
\end{thebibliography}

\GOneEnd
